\documentclass[12pt,a4paper,twoside]{book}
\usepackage{fontspec}
\usepackage{xgreek}

%Long Tables
\usepackage{longtable}

%Table of Contents
\usepackage{tocloft}

%Άδεια creative commons
\usepackage[scale=1]{ccicons}


%Ρυθμίσεις για συνδέσμους
\usepackage[hyphens]{url}
\usepackage[hidelinks]{hyperref}
%Χρώματα
\usepackage[dvipsnames*]{xcolor}
%Μαθηματικές εκφράσεις
%\usepackage{unicode-math}
%\setmathfont{Asana Math}

%Quotes package
\usepackage{epigraph}
\setlength\epigraphwidth{8cm}

%Κώδικας γλωσσών προγραμματισμού
\usepackage{listings}

%Κενό μεταξύ αντικειμένων λίστας
\usepackage{enumitem}
\setlist{noitemsep}

%Fancy Headers package
\usepackage{fancyhdr} 
\pagestyle{fancy}
\fancyhf{}
\fancyhead[LE,RO]{\nouppercase\rightmark}
\fancyhead[RE,LO]{\nouppercase\leftmark}
\fancyfoot[CE,CO]{\thepage}
\fancyfoot[LE,RO]{}

%Font Awesome Icons
\usepackage{fontawesome}


\usepackage{hologo}
\usepackage[top=2.54cm, bottom=2.54cm, left=3.17cm, right=3.17cm]{geometry}
\usepackage[Greek,Latin]{ucharclasses}
\usepackage{csquotes}
\usepackage{enumitem}
%Γραφικά
\usepackage{graphicx}
\graphicspath{{./images/}}
\usepackage{wrapfig}
%Διάταξη τίτλων
\usepackage{titlesec}
  
%Στοιχεία PDF
\author{Γεώργιος Σοφιανός}
\title{GPGPU Programming}
\setmainfont[Kerning=On,Mapping=tex-text]{GFS Artemisia}
%Νέες εντολές
\newcommand{\HRule}{\rule{\linewidth}{0.5mm}}
%Βιβλιογραφία
\usepackage[
  backref=true,
  backend=biber,
  sortcites, ]
  {biblatex}
\DefineBibliographyStrings{english}{%
    backrefpage  = {see p.}, % for single page number
    backrefpages = {see pp.} % for multiple page numbers
}
\addbibresource{bibliography.bib}


%Package to put image on absolute destinations
\usepackage[abs]{overpic}



%Fancy Sections - TODO LIST
\newcommand{\Bsection}[1]{%
%\noindent\hspace*{-11.6mm}
\begin{overpic}[width=\textwidth]%
	{banner.png}
	\put(0,0){\parbox{\textwidth}{\vspace{-20mm}\bfseries\color{white}\Large\centering#1}}
\end{overpic}
}

% New environment for Long Tables
\newenvironment{apptable}[2]
{\small
\begin{longtable}[c]{|p{3cm}|p{6cm}|l|l|}
\caption{#1\label{#2}}\\

\hline Όνομα & Περιγραφή & Επιτάχυνση & Multi-GPU \\ \hline
\endfirsthead
 
\hline
\multicolumn{4}{|c|}{Συνέχεια πίνακα \ref{#2}}\\
\hline
Όνομα & Περιγραφή & Επιτάχυνση & Multi-GPU \\ 
\hline
\endhead

\hline
\endfoot
 
\hline
\endlastfoot}
{\end{longtable}
\normalsize}

%Japanese font for Kanji
\newfontfamily\kanji{MingLiU}

%more fonts
\newfontfamily\arial{Arial}

%Chapter title configuration
%\titlespacing*{\chapter}{0pt}{-2\baselineskip}{1em}
\titleformat{\chapter}[display]
  {\Large\bfseries} % format
  {}                % label
  {0pt}             % sep
  {\huge\centering}           % before-code
  [
\vspace{-0.5ex}%
\rule{\textwidth}{0.3pt}
] % after-code
\titlespacing*{\chapter}{0cm}{-50pt}{0pt}[0pt]

%Section title Configuration
\titleformat{\section}
  {\large\bfseries} % format
  {\thesection}      % label
  {0.5ex}             % sep
  {\large}           % before-code
  [
\vspace{-0.5ex}%
\rule{\textwidth}{0.3pt}
] % after-code

%Subsection title Configuration
\titleformat{\subsection}
  {\normalsize\bfseries} % format
  {\thesubsection}                % label
  {0.5ex}             % sep
  {\normalsize}           % before-code
%Sections on new page
\newcommand{\sectionbreak}{\clearpage}

%Οριζόντια γραμμή μεγάλη
\newcommand{\HRule}{\rule{\linewidth}{0.5mm}}
\begin{document}
\setlist[itemize]{leftmargin=*}
\setlist[enumerate]{leftmargin=*}
\frontmatter
\pagenumbering{Roman}
\begin{titlepage}
\begin{center}

\includegraphics[scale=0.25]{Logo_teiath}~\\[1cm]

\textsc{\LARGE ΤΕΙ ΑΘΗΝΑΣ ΤΜΗΜΑ ΠΛΗΡΟΦΟΡΙΚΗΣ}\\[1.5cm]

\textsc{\Large Πτυχιακή Εργασία}\\[0.5cm]

% Τίτλος
\HRule \\[0.4cm]
{ \huge \bfseries Προγραμματισμός με τεχνικές GPGPU \\[0.4cm] }

\HRule \\[1.5cm]

% Φοιτητής και επόπτης
\begin{minipage}{0.4\textwidth}
\begin{flushleft} \large
\emph{Φοιτητής:}\\
\textsc{Γεώργιος Σοφιανός}\\
\emph{A.M:} 031053
\end{flushleft}
\end{minipage}
\begin{minipage}{0.4\textwidth}
\begin{flushright} \large
\emph{Επόπτης:} \\
\textsc{κ. Γεώργιος Μπαρδής}
\end{flushright}
\end{minipage}

\vfill

% Τέλος σελίδας
{\large \today}

\end{center}
\end{titlepage}
\chapter{Περίληψη}
\subsubsection{Περίληψη}
Σκοπός αυτής της εργασίας είναι η διερεύνηση και αξιολόγηση προτύπων και τεχνολογιών για προγραμματισμό γενικού σκοπού με χρήση μονάδων επεξεργασίας γραφικών, και η παρουσίαση εφαρμογών υψηλών υπολογιστικών απαιτήσεων οι οποίες εκμεταλλεύονται τις δυνατότητες αυτών των τεχνολογιών για επιτάχυνση και αύξηση των επιδόσεων τους.
\\
\\
\subsubsection{Abstract}
The purpose of this thesis is the research and evaluation of patterns and technologies used by general purpose programming on graphics processing units, and the demonstration of high processing computing applications which exploit the capabilities of these technologies in order to accelerate and increase their performance.

\newpage
\tableofcontents
\chapter{Εισαγωγή}
\epigraph{τεστ τεστ δοκιμή}{Γιώργος}
Ιστορία του GPGPU Programming:
Τα αρχικά GPGPU πηγάζουν απο την φράση General Purpose computation on Graphics Processing Units, ή αλλιώς γνωστή ώς GPU Computing, δηλαδή υπολογισμός γενικού σκοπού σε μονάδες επεξεργασίας γραφικών. Οι GPUs, είναι επεξεργαστές υψηλών επιδόσεων με δυνατότητα πολύ υψηλού υπολογισμού και διεκπεραιωτικότητας δεδομένων. Σχεδιασμένες αρχικά για γραφικά υπολογιστών με αρκετές δυσκολίες στον προγραμματισμό τους, οι σημερινές μονάδες επεξεργασίας γραφικών είναι παράλληλοι επεξεργαστές γενικής χρήσης με υποστήριξη για προσβάσιμες προγραμματιστικές διεπαφές και βιομηχανικά πρότυπα γλωσσών όπως η C. Οι προγραμματιστές που μεταφέρουν τις εφαρμογές τους σε GPUs συνήθως πετυχαίνουν ταχύτητες πολλαπλάσιες από ότι μια αντίστοιχη εφαρμογή ειδικά βελτιστοποιημένη για κεντρική μονάδα επεξεργασίας (CPU).
Ο όρος GPGPU δημιουργήθηκε από τον \href{http://en.wikipedia.org/wiki/Mark_Harris_(programmer)}{Mark Harris} το 2002 όταν συνειδητοποίησε ότι αναπτυσσόταν μια τάση για χρήση των μονάδων επεξεργασίας γραφικών για εφαρμογές που δεν είχαν σχέση με γραφικά. 

Από το 2012, οι GPU έχουν αναπτυχθεί σε συστήματα πολυπύρηνων επεξεργαστών παράλληλου υπολογισμού δίνοντας μας την δυνατότητα για πολύ αποδοτικό χειρισμό μεγάλου όγκου δεδομένων. Αυτός ο σχεδιασμός είναι πιο αποδοτικός από ότι οι κεντρικές μονάδες επεξεργασιάς (CPU) για αλγόριθμους όπου η επεξεργασία μεγάλου όγκου δεδομένων γίνεται παράλληλα, όπως σε αλγορίθμους sort μεγάλων λιστών, μετασχηματισμό κυμάτων δυο διαστάσεων, προσομοίωση βιολογικών δυναμικών.
\mainmatter
\chapter{Υλοποιήσεις}
\section{Ιστορία}
Παραδοσιακά, η σειρά αγωγών των γραφικών, αποτελείται από τις καταστάσεις μετατροπή και φωτισμός, συναρμολόγηση αρχέγονων, μετατροπή σε pixels, και σκίαση. Οι πρώτες GPU είχαν όλες τις λειτουργίες που χρειάζονται για να εκτελεστεί η σειρά αγωγών, αλλά με τον καιρό όλο και περισσότερες καταστάσεις έγιναν δυνατό να προγραμματιστούν με την έλευση ειδικών επεξεργαστών, όπως επεξεργαστές κορυφών και τεμάχια επεξεργαστών, που κατέστησαν κάποιες λειτουργίες πιο ευέλικτες.\\
Όταν οι τιμές συνέχισαν να πέφτουν ενώ η υπολογιστική δύναμη αυξανόταν, η ερευνητική κοινότητα σκέφτηκε τρόπους να αξιοποιηθεί αυτή η δύναμη για τον υπολογισμό δύσκολων λειτουργιών. Όμως, καθώς η δυνατότητα των επεξεργαστών ήταν περιορισμένη και η διεπαφή προγραμματιστικής διεπαφής (API) των οδηγών γραφικών ήταν σχεδιασμένη για να υλοποιεί συγκεκριμένα την σειρά αγωγών, έπρεπε να ληφθούν υπόψιν πολλές παράμετροι. Για παράδειγμα, όλα τα δεδομένα έπρεπε να κωδικοποιηθούν σε υφές ως πίνακες δυο διαστάσεων που αναπαριστούν pixel, με περιεχόμενο τιμές χρωμάτων και κάποιο κανάλι alpha για την διαφάνεια. Επιπλέον, οι υφές είναι αντικείμενα μόνο προσπελάσιμα, και δεν επαναγράφονταν, κάτι που ανάγκαζε τους προγραμματιστές να αποθηκεύουν κάθε φορά καινούρια υφή με τις αλλαγές. Τέλος, οι περισσότερες GPU υποστήριζαν μόνο λειτουργίες μονής κινητής υποδιαστολής, αναγκάζοντας τους προγραμματιστές να προσομοιώνουν λογικές λειτουργίες.\\
Αυτοί οι περιορισμοί, ήταν ο μεγαλύτερος λόγος που ώθησε τους κατασκευαστές GPU (AMD,NVIDIA,INTEL), να δημιουργήσουν προγραμματιστικές διεπαφές ειδικές για την κοινότητα του GPGPU και να εξελίξουν τις συσκευές τους για καλύτερη υποστήριξη.

Το πεδίο του προγραμματισμού γενικής χρήσης έχει αναπτυχθεί με ταχύτατους ρυθμούς τα τελευταία χρόνια, έτσι ώστε τώρα υπάρχουν αρκετές υλοποιήσεις για τον προγραμματισμό των μονάδων επεξεργασίας γραφικών. Πρόσφατα, έχουν γίνει προσπάθειες δημιουργίας προτύπων.
Ο προγραμματισμός των GPUs αναπτύχθηκε όταν το CUDA και το Stream κατέφθασαν στο τέλος του 2006. Αυτές οι διεπαφές και οι γλώσσες, σχεδιάστηκαν από τις εταιρίες κατασκευής των GPUs σε πολύ κοντινή σχέση με το υλικό, το οποίο αποτέλεσε μεγάλο βήμα προς ένα πιο εύχρηστο, ταιριαστό και μελλοντικά-ασφαλές προγραμματιστικό μοντέλο. Η ανοιχτή γλώσσα προγραμματισμού (OpenCL) δημιουργήθηκε για να παρέχει ένα γενικό API ετερογενή υπολογισμού σε διάφορες μορφές παράλληλων συσκευών, συμπεριλαμβανομένου μονάδων επεξεργασίας γραφικών, πολυπύρηνων κεντρικών μονάδων επεξεργασίας, κ.α
\section{CUDA}
\subsection{Εισαγωγή}
To CUDA είναι μια πλατφόρμα παράλληλου υπολογισμού, που δημιουργήθηκε από την NVIDIA και υλοποιήθηκε στις κάρτες γραφικών τις οποίες παράγει η ίδια. Το CUDA δίνει στους προγραμματιστές άμεση πρόσβαση στο σετ εικονικών εντολών και την μνήμη των στοιχείων του παράλληλου υπολογισμού σε κάρτες γραφικών NVIDIA. 

Αξιοποιώντας το CUDA, οι κάρτες γραφικών(GPU) μπορούν να χρησιμοποιηθούν για υπολογισμό γενικής χρήσης (δηλαδή όχι αποκλειστικά για γραφικά).Οι GPU έχουν μια αρχιτεκτονική παράλληλης εξόδου η οποία δίνει έμφαση στην εκτέλεση πολλών threads με μικρή ταχύτητα, σε αντίθεση με τις CPU όπου εκτελείται ένα thread με μεγάλη ταχύτητα. 

\begin{figure}[h]
\centering
\includegraphics[width=\linewidth]{nvidia-cuda}
\caption{Λογότυπο NVIDIA CUDA\cite{figure-3}}
\end{figure}

Η πλατφόρμα CUDA είναι προσβάσιμη στους προγραμματιστές μέσω βιβλιοθηκών, εντολών μεταγλώττισης, και προεκτάσεων σε γλώσσες προγραμματισμού βιομηχανικής κλίμακας, όπως η C, C++ και Fortran.

Οι προγραμματιστές της C/C++, χρησιμοποιούν το CUDA C/C++, μεταγλωττισμένο με το nvcc, έναν LLVM βασισμένο μεταγλωττιστή, και οι προγραμματιστές της Fortran χρησιμοποιούν το CUDA Fortran, μεταγλωττισμένο με τον μεταγλωττιστή PGI CUDA Fortran απο το The Portland Group. Εκτός από τα παραπάνω, η πλατφόρμα CUDA υποστηρίζει και άλλες διεπαφές υπολογισμού, όπως το OpenCL του Khronos Group, το DirectCompute της Microsoft, και το C++ AMP.

Στην βιομηχανία των υπολογιστών, οι GPUs δεν χρησιμοποιούνται μόνο για τα γραφικά αλλά και στους υπολογισμούς φυσικής παιχνιδιών (π.χ καπνός, φωτιά, ροή υγρών). Γνωστά παραδείγματα αποτελούν οι μηχανές PhysX και η Bullet. Το CUDA επίσης χρησιμοποιείται για να επιταχύνει μη-γραφικές εφαρμογές στην βιοπληροφορική, στην κρυπτογραφία, και σε πολλά άλλα πεδία.\cite{cuda-1}

Γενικότερα, η υπολογιστική δύναμη της GPU, βασίζεται στην παράλληλη αρχιτεκτονική της. Για αυτό, η πλατφόρμα του CUDA παρουσιάζει το νήμα(thread) ως το μικρότερο στοιχείο παραλληλισμού. Όμως, σε σύγκριση με την κεντρική μονάδα επεξεργασίας, τα νήματα της GPU έχουν μικρότερο κόστος χρήσης πόρων και μικρότερο κόστος δημιουργίας και αντικατάστασης.

Σημειώνεται ότι οι GPU είναι αποτελεσματικές, μόνο όταν τρέχει μεγάλος αριθμός απο τέτοια νήματα. Μια ομάδα από νήματα, που εκτελούνται παράλληλα, επικοινωνούν και συγχρονίζονται μεταξύ τους ονομάζεται block. Ο μέγιστος αριθμός των νημάτων σε ενα block είναι ένας περιορισμός που υπάρχει στην κάθε μονάδα γραφικής επεξεργασίας. Τέλος, μια ομάδα από blocks τα οποία έχουν την ίδια διάσταση και εκτελούνται απο το ίδιο πρόγραμμα CUDA παράλληλα, ονομάζεται πλέγμα.

Για να επιτρέψει βέλτιστη επίδοση για διαφορετικά πρότυπα, το CUDA εκτελεί ένα ιεραρχικό μοντέλο μνήμης, αντίθετα με τα παραδοσιακά μοντέλα που συναντάμε συνήθως στους υπολογιστές. Ο υπολογιστής και η συσκευή, έχουν τις δικές τους περιοχές μνήμης, τις οποίες ονομάζουν host memory και device memory, αντίστοιχα. Το CUDA παρέχει βελτιστοποιημένες λειτουργίες για να μεταφέρει δεδομένα από και προς αυτούς τους ξεχωριστούς χώρους.
Κάθε νήμα κατέχει το δικό του αρχείο καταχώρησης, το οποίο μπορεί να προσπελαστεί και να εγγραφεί.

Επιπλέον, μπορεί να προσπελάσει το δικό του αντίγραφο της τοπικής μνήμης. Όλα τα νήματα στο ίδιο πλέγμα μπορούν να προσπελάσουν και να γράψουν στην περιοχή της κοινόχρηστης μνήμης (shared memory). Για να αποφευχθούν κίνδυνοι από ταυτόχρονη προσπέλαση, μηχανισμοί συγχρονισμού νημάτων πρέπει να χρησιμοποιηθούν. Η κοινόχρηστη μνήμη, είναι οργανωμένη σε ομάδες που ονομάζονται τράπεζες, οι οποίες μπορούν να προσπελαστούν παράλληλα. Όλα τα νήματα έχουν επίσης πρόσβαση στον χώρο μνήμης που ονομάζεται καθολική μνήμη (global memory) και στις περιοχές που ονομάζονται μνήμη σταθερών (constant memory) και μνήμη υφής (texture memory).\cite{cuda-2}

\subsection{Πλεονεκτήματα}
Το CUDA έχει τα εξής πλεονεκτήματα σε σχέση με τους παραδοσιακούς τρόπους υπολογισμού γενικής χρήσης που εκτελούνται μέσω προγραμματιστικών διεπαφών γραφικών:
\begin{itemize}
\item Διασκορπισμένες προσπελάσεις - ο κώδικας μπορεί να διαβαστεί από αυθαίρετες διευθύνσεις στην μνήμη.
\item Ενοποιημένη εικονική μνήμη (CUDA 6)
\item Κοινόχρηστη μνήμη - το CUDA εκθέτει μια γρήγορη περιοχή κοινόχρηστης μνήμης (μέχρι 48KB για κάθε επεξεργαστή) η οποία μπορεί να μοιραστεί ανάμεσα στα threads. Αυτή μπορεί να χρησιμοποιηθεί σαν κρυφή μνήμη διαχειρίσιμη απο τον χρήστη, επιτρέποντας μεγαλύτερο εύρος δεδομένων απο ότι είναι δυνατό με τις προσπελάσεις υφών.
\item Πιο γρήγορες μεταφορτώσεις και προσπελάσεις από και προς την GPU
\item Πλήρης υποστήριξη για ακέραιες και bitwise λειτουργίες, για παράδειγμα τις προσπελάσεις υφών.
\end{itemize}
\subsection{Περιορισμοί}
 Το CUDA δεν υποστηρίζει ολόκληρο το πρότυπο της γλώσσας C, καθώς τρέχει μέσω ενός μεταγλωττιστή C++, ο οποίος εμποδίζει συγκεκριμένα μέρη της γλώσσας C να μεταγλωττιστούν.

\subsection{Μνήμη}
\subsubsection{Κοινή μνήμη και συγχρονισμός}
Ο μεταγλωττιστής CUDA C μεταχειρίζεται τις μεταβλητές στην κοινή μνήμη με διαφορετικό τρόπο απο ότι τις τυπικές μεταβλητές. Δημιουργεί ένα αντίγραφο για κάθε block που εκτελείται στην CPU. Κάθε νήμα σε αυτό το block μοιράζεται την μνήμη, αλλά τα νήματα δεν μπορούν να προσπελάσουν και να επεξεργαστούν το αντίγραφο της μεταβλητής που φαίνεται στα άλλα blocks.

Αυτό παρέχει ένα καλό τρόπο με τον οποίο τα νήματα μέσα σε ένα block μπορούν να επικοινωνούν και να συνεργάζονται στους υπολογισμούς. Επιπλέον, τα buffers κοινής μνήμης βρίσκονται πάνω στην GPU, με αυτον τον τρόπο, η καθυστέρηση στην προσπέλαση της κοινής μνήμης είναι πολύ μικρότερη από ότι στα τυπικά buffers, καθιστώντας την κοινή μνήμη πολύ αποδοτική.\cite{cuda-3}

Η επικοινωνία μεταξύ των νημάτων, αν και πολύ ενδιαφέρουσα, χρειάζεται έναν μηχανισμό για τον συγχρονισμό της. Για παράδειγμα, αν το νήμα Α γράψει μια τιμή στην κοινή μνήμη και θέλουμε το νήμα Β να κάνει κάτι με αυτήν την τιμή, δεν μπορούμε να ξεκινήσουμε το νήμα Β έως ότου γνωρίζουμε ότι η εγγραφή από το νήμα Α ολοκληρώθηκε. Χωρίς τον συγχρονισμό, θα είχαμε έναν αγώνα δρόμου όπου το σωστό αποτέλεσμα της εκτέλεσης θα εξαρτάται από μη ντετερμινιστικά χαρακτηριστικά του υλικού.
\subsubsection{Σταθερή μνήμη}
Έχουμε αναλύσει το πως οι μοντέρνες GPUs είναι εφοδιασμένες με τεράστιες δυνατότητες υπολογιστικής δύναμης. Το υπολογιστικό πλεονέκτημα που έχουν οι μονάδες επεξεργασίας γραφικών βοήθησε στην ανάπτυξη του προγραμματισμού γενικού σκοπού. Με εκατοντάδες αριθμητικές μονάδες στην GPU, συνήθως ο περιορισμός δεν είναι η αριθμητική απόδοση, αλλά το εύρος ζώνης της μνήμης. Υπάρχουν τόσες πολλές ALUs στους επεξεργαστές γραφικών, όπου πολλές φορές δεν προλαβαίνουμε να μεταφέρουμε τα δεδομένα αρκετά γρήγορα ώστε να διατηρήσουμε έναν υψηλό ρυθμό υπολογισμού. 

Η γλώσσα CUDA παρέχει άλλον έναν τύπο μνήμης γνωστή ως σταθερή μνήμη. Όπως φαίνεται από το όνομα, χρησιμοποιούμε την σταθερή μνήμη για δεδομένα που δεν αλλάζουν κατά την διάρκεια μιας εκτέλεσης πυρήνα. Το υλικό NVIDIA παρέχει 64KB σταθερής μνήμης που χρησιμοποιεί με διαφορετικό τρόπο από ότι την γενική μνήμη. Σε μερικές περιπτώσεις, η χρήση της σταθερής μνήμης αντί της γενικής μνήμης μειώνει το εύρος ζώνης της μνήμης. 

Δηλώνοντας την μνήμη σαν σταθερή, περιορίζουμε την χρήση της σε ανάγνωσης μόνο. Λόγω αυτού του περιορισμού, περιμένουμε να κερδίσουμε κάτι από αυτήν την διαδικασία. Η χρήση σταθερής μνήμης μπορεί να μας διαφυλάξει εύρος μνήμης σε σχέση με την ανάγνωση των δεδομένων από την γενική μνήμη. Υπάρχουν δύο λόγοι γιατί η ανάγνωση από την σταθερή μνήμη των 64KB μπορεί να διαφυλάξει εύρος μνήμης:
\begin{itemize}
\item Μια ανάγνωση από την σταθερή μνήμη μπορεί να αναμεταδοθεί σε κοντινά νήματα, σώζοντας μας έως και 15 αναγνώσεις.
\item Η σταθερή μνήμη είναι cached, έτσι επόμενες αναγνώσεις της ίδιας διεύθυνσης δεν θα δημιουργήσουν επιπλέον κίνηση στην μνήμη.
\end{itemize}
Για να εξηγήσουμε αυτήν την δήλωση και το τι είναι τα κοντινά νήματα, σκεφτόμαστε την περίπτωση της ύφανσης. Στην ύφανση, το στημόνι (warp) αναφέρεται σε ένα σύνολο από νήματα, που υφαίνονται μαζί σε ένα ύφασμα. Στην αρχιτεκτονική CUDA, το στημόνι αναφέρεται σε μια συλλογή από 32 νήματα τα οποία υφαίνονται μεταξύ τους, και εκτελούνται αμφίδρομα. Σε κάθε γραμμή του κώδικα, κάθε νήμα από το στημόνι εκτελεί την ίδια διαδικασία, σε διαφορετικά δεδομένα.
Όσον αφορά τον χειρισμό της σταθερής μνήμης, το υλικό NVIDIA μπορεί να αναμεταδώσει μια απλή ανάγνωση σε μια ομάδα απο 16 νήματα: τα μισά από τα 32 νήματα που βρίσκονται στο στημόνι. Αν κάθε νήμα ζητάει δεδομένα απο την ίδια διεύθυνση της σταθερής μνήμης, η μονάδα επεξεργασίας γραφικών θα εκτελέσει μόνο ένα αίτημα ανάγνωσης και θα αναμεταδώσει τα δεδομένα σε όλα τα νήματα. Αν η ανάγνωση γίνεται απο την σταθερή μνήμη για μεγάλο όγκο πληροφοριών, θα δημιουργηθεί κίνηση που αντιστοιχεί μόνο στο 1/16 της κίνησης μνήμης που θα χρειαζόταν για την γενική μνήμη.

Τα πλεονεκτήματα όμως δεν σταματάνε εκεί. Επειδή έχουμε αποφασίσει να αφήσουμε άθικτη την μνήμη, το υλικό μπορεί να αποθηκεύσει τα σταθερά δεδομένα στην GPU. Έτσι μετά από την πρώτη ανάγνωση από την διεύθυνση της σταθερής μνήμης, όλες οι επόμενες αναγνώσεις δεν θα δημιουργήσουν επιπλέον κίνηση στην μνήμη. Δυστυχώς, μπορεί να υπάρξουν στιγμές που η απόδοση να μειώνεται λόγω της σταθερής μνήμης. Η αναμετάδοση μέρους από το στημόνι (warp) είναι ένα δίκοπο μαχαίρι. Αν και μπορεί να επιταχύνει την απόδοση όταν τα 16 νήματα διαβάζουν την ίδια διεύθυνση, μπορεί αντίστοιχα να μειώσει την απόδοση όταν τα 16 νήματα διαβάζουν διαφορετικές διευθύνσεις. Για παράδειγμα, αν τα 16 νήματα χρειάζονται διαφορετικά δεδομένα απο την σταθερή μνήμη, οι 16 αναγνώσεις θα γίνονταν με την σειρά, και θα χρειάζονταν 16 φορές περισσότερο χρόνο για να εκτελέσουν την εντολή. Αν διάβαζαν απο την γενική μνήμη, η εντολή θα δινόταν άμεσα. Σε αυτήν την περίπτωση, η ανάγνωση απο την σταθερή μνήμη θα ήταν πιο αργή από το να γινόταν ανάγνωση της γενικής μνήμης.\cite{cuda-4}


\subsection{CUDA streams}
Εξηγήσαμε το πως ο παράλληλος προγραμματισμός δεδομένων σε μια GPU μπορεί να δώσει εντυπωσιακά αποτελέσματα σε σχέση με την εκτέλεση σε CPU. Όμως υπάρχει ακόμα ένας τύπος παράλληλου υπολογισμού που μπορούμε να εκμεταλλευτούμε σε μια μονάδα επεξεργασίας γραφικών NVIDIA. Ο παραλληλισμός αυτός μοιάζει με αυτόν που συμβαίνει στα πολυ-νηματικά προγράμματα CPU. Αντί να εκτελούμε την ίδια διαδικασία σε πολλά στοιχεία δεδομένων όπως στον παραλληλισμό δεδομένων, ο παραλληλισμός έργων περιλαμβάνει παράλληλη εκτέλεση δύο η περισσότερων έργων. 

Σαν έργο μπορούμε να θεωρήσουμε μεγάλο αριθμό πραγμάτων. Για παράδειγμα, μια εφαρμογή μπορεί να εκτελεί δύο έργα: επανασχεδίαση του γραφικού περιβάλλοντος με ένα νήμα, και μεταφόρτωση μιας αναβάθμισης μέσω δικτύου με κάποιο άλλο νήμα. Αυτά τα έργα εκτελούνται παράλληλα, και ας μην έχουν τίποτε κοινό. Αν και ο παραλληλισμός έργων δεν είναι τόσο ευέλικτος όσο στις CPUs, συνεχίζει να μας προσφέρει ευκαιρίες για να αποκτήσουμε περισσότερη ταχύτητα από τις βασισμένες σε GPU εφαρμογές μας.

Τα CUDA streams, παίζουν μεγάλο ρόλο στην επιτάχυνση των εφαρμογών μας. Ένα CUDA stream αντιπροσωπεύει μια λίστα από GPU διεργασίες που θα εκτελεσθούν με συγκεκριμένη σειρά. Οι διεργασίες μπορούν να περιλαμβάνουν εκτελέσεις πυρήνων, αντιγραφές μνήμης, και εκκινήσεις/τερματισμούς συμβάντων ενός stream. H σειρά εισαγωγής των διεργασιών ορίζει την σειρά εκτέλεσης τους, με ευκαιρίες για αυτά τα έργα να εκτελεστούν παράλληλα.

%TODO: STUFF MISSING here 
\section{OpenCL}
\subsection{Εισαγωγή}
Το πρώτο GPGPU framework δημιουργήθηκε απο την NVIDIA και ήταν το CUDA. Το CUDA παρείχε στους χρήστες ένα προγραμματιστικό περιβάλλον σε C like γλώσσα για την GPU. Όμως ήταν κλειστού κώδικα και μπορεί να τρέχει μόνο σε NVIDIA κάρτες γραφικών. Λόγω της μεγάλης δημοτικότητας του CUDA, η ανάγκη για ένα ανοιχτό πρότυπο αρχιτεκτονικής που θα υποστηρίζει διάφορα είδη συσκευών απο διάφορους κατασκευαστές γινόταν όλο και πιο σημαντική. Έτσι τον Ιούνιο του 2008 το Khronos Group δημιούργησε το OpenCL 1.0. Αρκετοί κατασκευαστές σταδιακά παρείχαν εργαλεία για προγραμματισμό σε OpenCL συμπεριλαμβανομένων των: Nvidia OpenCL Drivers and Tools, AMD APP SDK,Intel SDK for OpenCL applications,IBM Server with OpenCL development Kit, κ.α. Σήμερα το OpenCL επιτρέπει πολυπύρηνο προγραμματισμό, προγραμματισμό GPU, κ.α.\cite{opencl-1}   

Το OpenCL είναι ένα πλαίσιο για κατασκευή εφαρμογών που εκτελούνται σε ετερογενή συστήματα που αποτελούνται από κεντρικές μονάδες επεξεργασίας(CPU), μονάδες επεξεργασίας γραφικών(GPU), επεξεργαστές ψηφιακών σημάτων(DPS), συστοιχίες προγραμματιζόμενων θυρίδων(FPGA), και άλλους επεξεργαστές. Το OpenCL περιέχει μια γλώσσα, υποσύνολο του ISO C99 με επεκτάσεις, για τον προγραμματισμό αυτών των συσκευών, προγραμματιστικές διεπαφές εφαρμογών(API) για τον έλεγχο της πλατφόρμας και την εκτέλεση προγραμμάτων στις υπολογιστικές συσκευές. Το OpenCL παρέχει παράλληλο υπολογισμό χρησιμοποιώντας παραλληλισμό διεργασιών και δεδομένων.\\
Αποτελεί το πρώτο ανοιχτό, ελεύθερο από τέλη αδειών πρότυπο για cross-platform, παράλληλο προγραμματισμό μοντέρνων επεξεργαστών, που χρησιμοποιούνται συνήθως σε προσωπικούς υπολογιστές, διακομιστές, και φορητές/ενσωματωμένες συσκευές. Το OpenCL (Open Computing Language) βελτιώνει αισθητά την ταχύτητα και την απόκριση μεγάλου εύρους εφαρμογών σε διάφορες κατηγορίες αγορών από παιχνίδια και ψυχαγωγία, μέχρι επιστημονικές εφαρμογές και εφαρμογές υγείας. Συντηρείται απο το μή-κερδοσκοπικό τεχνολογικό συνεταιρισμό Khronos Group. Έχει υιοθετηθεί από πολλές μεγάλες εταιρίες όπως η Apple, Intel, Qualcomm, AMD, Nvidia,Samsung,ARM Holdings.
\begin{figure}[h]
	\includegraphics[width=0.5\linewidth]{opencl_logo}
	\centering
	\caption{Λογότυπο OpenCL}
\end{figure}
\subsection{Επισκόπηση}
Το OpenCL ορίζει μια διεπαφή προγραμματισμού εφαρμογών με την οποία επιτρέπει στα προγράμματα που τρέχουν στον οικοδεσπότη να εκτελέσουν πυρήνες στην συσκευή υπολογισμού, και να διαχειριστούν την μνήμη της συσκευής, που είναι ξεχωριστή από την μνήμη του οικοδεσπότη. Τα προγράμματα του OpenCL είναι σχεδιασμένα ώστε να μεταγλωττίζονται την ώρα της εκτέλεσης, και με αυτόν τον τρόπο γίνεται δυνατό να εκτελεστούν σε διάφορες συσκευές. Το πρότυπο του OpenCL ορίζει διεπαφές προγραμματισμού για γλώσσα C και C++. Διεπαφές επίσης υπάρχουν και για άλλες γλώσσες, όπως Python, Julia, και Java. Μια εφαρμογή του OpenCL προτύπου αποτελείται απο μια βιβλιοθήκη που υλοποιεί την διεπαφή για C και C++, και έναν μεταγλωττιστή OpenCL για τις συσκευές υπολογισμού.\\
\subsubsection{Ιεραρχία μνήμης}
Το OpenCL ορίζει ιεραρχία τεσσάρων επιπέδων για την μνήμη των συσκευών υπολογισμού:
\begin{itemize}
\item Καθολική μνήμη: διαμοιράζεται σε όλες τις συσκευές υπολογισμού, αλλά έχει μεγάλη καθυστέρηση απόκρισης
\item Μνήμη προσπέλασης: μικρότερη, χαμηλή καθυστέρηση απόκρισης, εγγράψιμη από την κεντρική μονάδα επεξεργασίας του οικοδεσπότη, αλλά όχι των συσκευών υπολογισμού.
\item Τοπική μνήμη: διαμοιράζεται σε πολλά στοιχεία υπολογισμού μιας συσκευής
\item Ιδιωτική μνήμη στοιχείου (καταχωρητές) 
\end{itemize}   
Δεν είναι απαραίτητο για όλες τις συσκευές να υλοποιήσουν την ιεραρχία της μνήμης στο υλικό. Η συνέπεια στα διάφορα επίπεδα της ιεραρχίας είναι χαλαρή, και επιβάλλεται μόνο από κατηγορηματικά στοιχεία συγχρονισμού, όπως τα εμπόδια.
\subsection{Ιστορία}
Το OpenCL δημιουργήθηκε αρχικά απο την Apple Inc., η οποία κατέχει τα πνευματικά δικαιώματα, και εξευγενίστηκε σε αρχική πρόταση σε συνεργασία με τεχνικές ομάδες της AMD,IBM,Qualcomm,Intel, και Nvidia. Η Apple καταχώρησε την πρόταση στο Khronos Group, και τον Ιούνιο του 2008 διαμορφώθηκε το Khronos Compute Working Group με αντιπρόσωπους από εταιρίες επεξεργαστών, μονάδων υλικού γραφικών, ενσωματωμένων-επεξεργαστών, και λογισμικού. Αυτό το group εργάστηκε για 5 μήνες ώστε να φέρει σε πέρας τον πρώτο προσδιορισμό για το OpenCL 1.0, ο οποίος κυκλοφόρησε τον Δεκέμβριο του 2008. 
\begin{table}[h]
	\begin{tabular}{|l|p{.9\linewidth}|}
	\hline
OpenCL 1.0 & Η AMD, αν και αρχικά εργαζόταν πάνω στο πρότυπο Close to Metal, αποφάσισε να στραφεί και να υποστηρίξει το OpenCL. Η Nvidia ανακοίνωσε πλήρης υποστήριξη στην εργαλειοθήκη υπολογισμού GPU. Το 2009, η IBM κυκλοφόρησε την πρώτη έκδοση του μεταγλωττιστή της με υποστήριξη για OpenCL\\ \hline
OpenCL 1.1 & Το OpenCL 1.1 επικυρώθηκε από το Khronos Group τον Ιούνιο του 2010, και προσθέτει σημαντικές λειτουργίες για βελτιωμένη ευελιξία παράλληλου προγραμματισμού, λειτουργικότητα, και επιδόσεις.\\ \hline
OpenCL 1.2 & Το OpenCL 1.2 ανακοινώθηκε τον Νοέμβριο του 2011 απο το Khronos Group, το οποίο προσθέτει αρκετές λειτουργίες σε σχέση με τις προηγούμενες εκδόσεις όσον αφορά τις επιδόσεις και χαρακτηριστικά για παράλληλο προγραμματισμό.\\ \hline
OpenCL 2.0 & Το OpenCL 2.0 επικυρώθηκε και κυκλοφόρησε τον Νοέμβριο του 2013 απο το Khronos Group και αποτελεί την τελευταία έκδοση του OpenCL. \\ \hline
	\end{tabular}
\end{table}

\subsection{Στόχοι}
Ο στόχος του OpenCL είναι να κάνει ορισμένους τύπους παράλληλου προγραμματισμού πιο εύκολους, και να παρέχει ανεξαρτήτου κατασκευαστή παράλληλη εκτέλεση κώδικα μέσω επιτάχυνσης υλικού. Το OpenCL είναι το πρώτο ανοιχτό, ελεύθερο πρότυπο για παράλληλο προγραμματισμό γενικού σκοπού ετερογενών συστημάτων. Παρέχει ένα προγραμματιστικό περιβάλλον που βοηθάει τους προγραμματιστές να γράψουν αποδοτικό, φορητό κώδικα για συστήματα υψηλής απόδοσης, προσωπικούς υπολογιστές, και κινητές συσκευές χρησιμοποιώντας ένα μείγμα πολυπύρηνων CPUs,GPUs, και DSPs.\\
Το OpenCL παρέχει στους προγραμματιστές ένα κοινό σετ εργαλείων εύκολης χρήσης, ώστε αυτοί να εκμεταλλευτούν οποιαδήποτε συσκευή που περιέχει οδηγό OpenCL για την εκτέλεση παράλληλου κώδικα. Το OpenCL framework ορίζει μια γλώσσα C like για την δημιουργία των πυρήνων, και ένα σετ απο APIs για την δημιουργία και την διαχείριση αυτών των πυρήνων. Οι πυρήνες είναι διαδικασίες οι οποίες μπορούν να εκτελούνται σε διαφορετικές συσκευές. Οι πυρήνες μεταγλωττίζονται απο έναν μεταγλωτιστή runtime, μέσω κατάλληλου προγράμματος. Αυτό επιτρέπει στα προγράμματα να εκμεταλλεύονται όλες τις συσκευές ενός συστήματος με ένα σετ φορητών υπολογιστικών πυρήνων.
\subsection{Μοντέλο εκτέλεσης}
Τα κύρια μέρη εκτέλεσης ενός προγράμματος OpenCL είναι ο πυρήνας και το πρόγραμμα ξενιστή. Οι πυρήνες εκτελούνται στην συσκευή OpenCL και το πρόγραμμα ξενιστή, στον υπολογιστή που εκτελείται το πρόγραμμα. Ο σκοπός του προγράμματος ξενιστή είναι να δημιουργήσει και να ζητήσει την πλατφόρμα και τις ιδιότητες της συσκευής, να ορίσει το περιεχόμενο, να κατασκευάσει τον πυρήνα, και να διαχειριστεί την εκτέλεση των πυρήνων. Όταν καταχωρηθεί ο πυρήνας απο τον ξενιστή στην συσκευή, δημιουργείται ένα Ν διαστάσεων χώρος ευρετηρίου, με το Ν να είναι από 1 έως 3. Κάθε περιστατικό πυρήνα δημιουργείται στις συντεταγμένες του χώρου ευρετηρίου. Αυτό το περιστατικό ονομάζεται αντικείμενο εργασίας και ο χώρος ευρετηρίου καλείται NDRange.

\subsection{WebCL}
Με την δημοτικότητα των εφαρμογών web να αυξάνεται, υπάρχει ανάγκη για αύξηση της απόδοσης παράλληλης υπολογιστικής επεξεργασίας για να επιταχύνουμε τις διεργασίες εντατικού υπολογισμού σε εφαρμογές ιστού. Αυτές οι εφαρμογές περιλαμβάνουν για παράδειγμα, παρουσίαση απεικόνισης δεδομένων, επεξεργασία εικόνας και βίντεο, παιχνίδια τριών διαστάσεων, υπολογιστική φωτογραφία, κρυπτογραφία. Παρέχοντας στους προγραμματιστές με ένα πρότυπο Javascript API και μια γλώσσα φορητού προγραμματισμού, το WebCL επιτρέπει παράλληλο υπολογισμό σε ετερογενή πολυ-πύρηνα συστήματα σε μια σωρεία συσκευών, συμπεριλαμβανομένου φορητών, σταθερών, και διακομιστών. 

Το WebCL 1.0 ορίζει ένα Javascript binding στο πρότυπο OpenCL για ετερογενή παράλληλο υπολογισμό, ενώ επιτρέπει σε εφαρμογές ιστού να εκμεταλλευτούν τις δυνατότητες της GPU και τον παράλληλο υπολογισμό πολυπύρηνων CPU, μέσα απο έναν Web Browser, ενεργοποιώντας σημαντική επιτάχυνση των εφαρμογών όπως επεξεργασία βίντεο και εικόνας, και ανώτερης εξομοίωσης φυσικής για παιχνίδια WebGL. Το WebCL έχει αναπτυχτεί σε στενή συνεργασία με την κοινότητα του web, και παρέχει την δυνατότητα να επεκταθούν οι δυνατότητες των HTML5 browsers ώστε να επιταχύνουν τις εφαρμογές υψηλών απαιτήσεων υπολογισμού και πλούσιου οπτικού υπολογισμού. Το WebCL έχει σχεδιαστεί και συντηρείται από το μη κερδοσκοπικό ίδρυμα Khronos Group. Οι ολοκληρωμένες προδιαγραφές της πρώτης έκδοσης του WebCL ανακοινώθηκαν τον Μάρτιο του 2014.

Ανάλογα την υλοποίηση, Το WebCL επιτρέπει στις διεργασίες να εκτελούνται ταυτόχρονα με την Javascript. Συγκεκριμένα, είναι δυνατόν για την εφαρμογή να επεξεργαστεί ένα buffer καθώς αντιγράφεται από ένα αντικείμενο WebCLMemoryObject. Για να αποτρέψουμε την φθορά των δεδομένων, οι εφαρμογές θα πρέπει να μην επεξεργάζονται τα buffers που έχουν σημαδευτεί για ασύγχρονη ανάγνωση/εγγραφή, έως ότου η σχετική ουρά εντολών WebCL έχει τελειώσει την εκτέλεση της.

\begin{figure}[h]
	\includegraphics[scale=1]{webcl-logo}
	\centering
	\caption{Λογότυπο WebCL}
\end{figure}
\begin{itemize}
\item Khronos Launching new WebCL initiative
	\begin{itemize}
	\item Ανακοινώθηκε τον Μάρτιο του 2011
	\item API definitions already underway
	\end{itemize}
\item Javascript binding για OpenCL
	\begin{itemize}
	\item Η ασφάλεια πρώτη προτεραιότητα
	\end{itemize}
\item Πολλές περιπτώσεις χρήσης
	\begin{itemize}
	\item Μηχανές φυσικής για συμπλήρωση του WebGL
	\item Επεξεργασία εικόνας και βίντεο σε browser
	\end{itemize}
\item Πολύ στενή σχέση με το πρότυπο OpenCL
	\begin{itemize}
	\item Maximum flexibility
	\item Foundation for higher-level middleware
	\end{itemize}
\end{itemize}


\subsubsection{Υλοποιήσεις}
Αυτήν την στιγμή, μόνο μια δοκιμαστική έκδοση του Chromium browser υποστηρίζει το WebCL, γιατί η τεχνολογία είναι σχετικά καινούρια. Όμως το webCL μπορεί να εκτελεστεί και ως ένα πρόσθετο. Για παράδειγμα, η Nokia έχει δημιουργήσει ένα πρόσθετο WebCL για τον Mozilla Firefox. Η Mozilla έχει ανακοινώσει ότι δεν θα υποστηρίξει το WebCL, λόγω των Compute Shader του OpenGL ES 3.1. Σαν παράδειγμα αναφέρουμε υλοποιήσεις που προσπαθούν να ακολουθήσουν τις προδιαγραφές WebCL 1.0
\begin{itemize}
\item Samsung (WebKit) - https://github.com/SRA-SiliconValley/webkit-webcl
\item Nokia (Firefox) - http://webcl.nokiaresearch.com/
\item Motorola Mobility node-webcl (node.js) - https://github.com/Motorola-Mobility/node-webcl
\end{itemize}


\subsubsection{Ασφάλεια}
Το WebCL έχει δημιουργηθεί δίνοντας μεγάλη έμφαση στην ασφάλεια. Μερικά από τα φαινόμενα που μπορεί να προκύψουν κατά την μεταγλώττιση ή και την εκτέλεση ενός πυρήνα WebCL είναι τα παρακάτω
\begin{itemize}
\item Πρόσβαση OOB(Out of Bounds) - Οι πυρήνες του WebCL πρέπει να μην επιτρέπουν πρόσβαση σε κρίσιμα μέρη της μνήμης, χωρίς να γίνονται διακρίσεις ανάλογα με τον τύπο της μεταβλητής.(private, global, constant). Αν ανιχνευθεί την ώρα της μεταγλώττισης, η πρόσβαση OOB πρέπει να παράγει ένα λάθος μεταγλώττισης. Την ώρα της εκτέλεσης, η προσπέλαση OOB πρέπει να επιστρέψει μηδέν, και οι εγγραφές να αγνοηθούν. Για λόγους ελέγχου του ορίου, η υλοποίηση μπορεί να μεταχειρίζεται τις μεταβλητές σαν ένα συνεχόμενο block μνήμης. Για τον έλεγχο των προσπελάσεων OOB, το Khronos Group ανέπτυξε τον WebCL Validator, ο οποίος αναγκάζει την αρχικοποίηση τοπικής μνήμης.
\item Αρχικοποίηση μνήμης - Για να σιγουρέψουμε ότι οι εφαρμογές δεν μπορούν να προσπελάσουν παλαιότερα δεδομένα που έχουν μείνει στην μνήμη απο προηγούμενες εφαρμογές, η υλοποίηση WebCL πρέπει να αρχικοποιεί όλα τα buffers και μεταβλητές στο μηδέν πριν επιτρέψει την πρόσβαση στην εφαρμογή μας. Αυτή η ανάγκη υπάρχει ανεξάρτητα από τον τύπο της μεταβλητής και ανεξάρτητα την συσκευή στην οποία εκτελείται ο κώδικας. Όπου είναι δυνατόν, το πρόσθετο OpenCL 1.2 "cl{\_}khr{\_}initialize{\_}memory" επιτρέπει στις υλοποιήσεις WebCL να αρχικοποιήσουν την τοπική μνήμη αυτόματα πριν την εκτέλεση ενός πυρήνα.
\item Άρνηση υπηρεσίας - Πυρήνες που εκτελούνται για μεγάλο χρονικό διάστημα ή πυρήνες εντατικών υπολογισμών (ή άλλες εντολές στην ουρά εντολών) μπορεί να προκαλέσουν αστάθεια στο σύστημα λόγω της υπερκατανάλωσης πόρων. Δεν είναι εύκολο να γίνει έλεγχος του συγκεκριμένου προβλήματος στο επίπεδο του WebCL. Οι απαραίτητοι μηχανισμοί, όπως μετρητές watchdog και προληπτικός πολυ-νηματικός προγραμματισμός, πρέπει να παρέχονται από τον οδηγό OpenCL και το λειτουργικό σύστημα. Σε συστήματα όπου οι παραπάνω μηχανισμοί είναι διαθέσιμοι, οι υλοποιήσεις WebCL πρέπει να τους χρησιμοποιούν για να:
	\begin{enumerate}
		\item Ανιχνεύουν προβληματικούς πυρήνες και άλλες εντολές. Μια εντολή θεωρείται προβληματική όταν εκτελείται για πολύ μεγάλο χρονικό διάστημα, ή καταναλώνει υπερβολικά πολλούς πόρους συστήματος
		\item Τερματίζει τα περιεχόμενα που σχετίζονται με τις προβληματικές εντολές, πριν αυτές καταστήσουν την συσκευή OpenCL απαθής, και την αναγκάσουν να χρειαστεί επανεκκίνηση.
	\end{enumerate}
\end{itemize}
Όπου αυτό είναι δυνατόν, το πρόσθετο OpenCL 1.2 "cl{\_}kh{\_}-terminate{\_}context" μπορεί να χρησιμοποιηθεί για γρήγορο τερματισμό ενός περιεχομένου, αν για παράδειγμα κάποιος πυρήνας εκτελείται για πολύ μεγάλο χρονικό διάστημα ή ένα από τα προγράμματα τερματιστεί λόγω εξαιρέσεων.

\section{Compute Shaders}
Τα shaders υπολογισμού είναι μια κατάσταση shader που χρησιμοποιείται σχεδόν αποκλειστικά για υπολογισμούς αυθαίρετης πληροφορίας.Αν και μπορεί να χρησιμοποιηθεί για απόδοση, συνήθως χρησιμοποιείται για διεργασίες που δεν σχετίζονται άμεσα με σχεδιασμό τριγώνων και pixel.\cite{computeshaders-1}
\subsection{Shaders}
\subsubsection{Εισαγωγή}
Στον τομέα των γραφικών υπολογιστών, ένα shader είναι ένα πρόγραμμα που χρησιμοποιείται για να εκτελέσει την λεγόμενη σκίαση: την παραγωγή συγκεκριμένου επιπέδου χρώματος μέσα σε μια εικόνα, ή την παραγωγή ειδικών εφέ ή μετατροπές βίντεο. Ένας όρος που περιγράφει την σκίαση είναι "ένα πρόγραμμα που μαθαίνει τον υπολογιστή πώς να ζωγραφίσει κάτι με έναν ειδικό και μοναδικό τρόπο".\\
Τα shader υπολογίζουν αποδόσεις εφέ σε υλικό γραφικών, με ένα μεγάλο βαθμό ευλυγισίας. Τα περισσότερα shader είναι σχεδιασμένα για χρήση σε μονάδα επεξεργασίας γραφικών (GPU), όμως αυτό δεν είναι αποκλειστική ανάγκη. Οι γλώσσες σκίασης, χρησιμοποιούνται συνήθως για να προγραμματίσουν την γραμμή σωλήνα απόδοσης της GPU. Η θέση, η απόχρωση, ο κορεσμός, η φωτεινότητα, και η αντίθεση όλων των στοιχείων, κορυφών, ή υφών, χρησιμοποιούνται για να αποδώσουν μια τελική εικόνα που μπορούμε να επεξεργαστούμε απευθείας με χρήση αλγορίθμων ορισμένων στα shader, είτε με αλλαγές απο εξωτερικές μεταβλητές που εισάγει το πρόγραμμα το οποίο καλεί τον shader.
Τα shader χρησιμοποιούνται πολύ στην κινηματογραφική επεξεργασία, στις εικόνες που αποδίδονται απο τον υπολογιστή, αλλα και σε παιχνίδια υπολογιστών, για να παράγουν ένα μεγάλο αριθμό απο εφέ. Εκτός απο τα απλά μοντέλα φωτισμού, μερικά απο τα πολύπλοκα εφέ επεξεργάζονται την εικόνα και προσθέτουν blur, light bloom, volumetric lightning, normal mapping, bokeh, cel shading,posterization,bump mapping,distortion,chroma keying, edge detection, motion detection, κ.α\\
Η σύγχρονη χρήση των shader ξεκίνησε απο την Pixar, τον Μάιο του 1988. Όσο οι μονάδες επεξεργασίας γραφικών εξελίσσονταν, οι γνωστές βιβλιοθήκες γραφικών ξεκίνησαν να υποστηρίζουν τα shader. Οι πρώτες κάρτες γραφικών υποστήριζαν μόνο pixel shader, αλλά σύντομα ακολούθησε η εισαγωγή των vertex shader όταν οι προγραμματιστές κατάλαβαν τις δυνατότητες τους. Τα shader γεωμετρίας εισήχθηκαν μόλις με το Direct3D 10 και το OpenGL 3.2
\subsubsection{Τύποι}
Υπάρχουν διάφοροι τύποι shader που χρησιμοποιούνται γενικά. Ενώ οι παλιές κάρτες γραφικών είχαν ξεχωριστό τρόπο επεξεργασίας στοιχείων για κάθε τύπο shader, οι καινούριες έχουν ενωμένους shader που έχουν την δυνατότητα να εκτελούν οποιονδήποτε τύπο shader. Αυτό επιτρέπει στις κάρτες γραφικών να έχουν πιο αποδοτική χρήση της επεξεργαστικής τους δύναμης.
\begin{itemize}
\item Vertex shaders - Μετατρέπουν κάθε θέση τρισδιάστατη στον εικονικό χώρο σε δισδιάστατη. Μπορούν να επεξεργαστούν ιδιότητες όπως η θέση, το χρώμα, και συντεταγμένες υφής, αλλα δεν μπορούν να δημιουργήσουν καινούρια vertices.
\item Pixel shaders - Υπολογίζουν το χρώμα και άλλες ιδιότητες ενός τεμαχίου. Οι απλές μορφές τους αποδίδουν ένα pixel εικόνας, ενώ οι πιο πολύπλοκες μορφές αποδίδουν πολλά. Στα τρισδιάστατα γραφικά, ένας pixel shader δεν μπορεί να παράγει πολύπλοκα εφέ, γιατί επεξεργάζεται μόνο ένα τεμάχιο, χωρίς κάποια γνώση της γεωμετρίας της οθόνης. Μπορούν όμως να εφαρμοστούν σε εφέ δύο διαστάσεων και χρησιμοποιούνται για επεξεργασία υφών. Για παράδειγμα είναι ο μόνος τύπος shader που μπορεί να λειτουργήσει σαν φίλτρο για μια ροή βίντεο.
\item Geometry shaders - Τα shader γεωμετρίας είναι σχετικά καινούριος τύπος shader. Μπορεί να δημιουργήσει γραφικά αρχικά στοιχεία όπως γραμμές, τρίγωνα. Τα shader γεωμετρίας εκτελούνται μετά απο τα vertex shader. Τυπικές χρήσεις των shader γεωμετρίας συμπεριλαμβάνουν γεωμετρική ψηφίδωση, εξώθηση σκιώδους όγκου, και απόδοση σε χάρτη κύβου. Για παράδειγμα, σε μια καμπυλωτή γραμμή, τα δεδομένα των στοιχείων εισάγονται σαν είσοδος στον shader, και αυτός αναλαμβάνει να δημιουργήσει αυτόματα επιπλέον γραμμές που δίνουν πιο εύστοχη απεικόνιση της καμπύλης.
\item Tesselation shaders - Με την κυκλοφορία του OpenGL 4.0 και του Direct3D 11, ένας καινούριος τύπος shader έχει προστεθεί που ονομάζεται shader ψηφίδωσης. Επιτρέπει στα αντικείμενα κοντά στην κάμερα να έχουν καλύτερη λεπτομέρεια, ενώ τα αντικείμενα που είναι πιο μακριά να έχουν μικρότερη λεπτομέρεια αλλά να μοιάζουν όμοια στην ποιότητα.
\item Compute shaders - Είναι ένας τύπος shader που χρησιμοποιείται αποκλειστικά για υπολογισμό αυθαίρετης πληροφορίας. Αν και μπορεί να χρησιμοποιηθεί για απόδοση, συνήθως χρησιμοποιείται για διεργασίες που δεν σχετίζονται άμεσα με σχεδιασμό τριγώνων και pixel.
\end{itemize}
\subsection{Εισαγωγή}
Τα shaders υπολογισμού λειτουργούν διαφορετικά από τις άλλες καταστάσεις shader. Όλες οι καταστάσεις shader έχουν προκαθορισμένου τύπου τιμές εισόδου, μερικές ενσωματωμένες και μερικές καθορισμένες από τον χρήστη. Η συχότητα στην οποία εκτελείται μια κατάσταση shader εξαρτάται από την φύση της κατάστασης. Για παράδειγμα, τα shader κορυφών εκτελούνται μια φορά για κάθε κορυφή. \\
Τα shader υπολογισμού λειτουργούν πολύ διαφορετικά. Ο "χώρος" στον οποίο ένα shader υπολογισμού λειτουργεί είναι αφηρημένος. Είναι στην κρίση του κάθε shader υπολογισμού να αποφασίσει τι σημαίνει αυτός ο "χώρος". Ο αριθμός των εκτελέσεων των shader υπολογισμού ορίζεται από την διεργασία που χρησιμοποιείται για να εκτελεστεί η υπολογιστική λειτουργία. Πιο σημαντικό από όλα, τα shader υπολογισμού δεν έχουν εισόδους καθορισμένες απο τον χρήστη και ούτε καμία έξοδο. Οι ενσωματωμένες είσοδοι ορίζουν μόνο το πού στον "χώρο" της εκτέλεσης βρίσκεται ένας συγκεκριμένος shader υπολογισμού.\\
Έτσι, αν κάποιος shader υπολογισμού πρέπει να πάρει κάποιες τιμές σαν είσοδο, είναι στην ευθύνη του shader να αποκτήσει τα δεδομένα, μέσω πρόσβασης υφών, αυθαίρετης φόρτωσης εικόνας, ή άλλες μορφές διεπαφής. Παρομοίως, αν ένας shader υπολογισμού υπολογίζει κάτι, θα πρέπει να το αποθηκεύσει σε μια εικόνα ή σε ένα block αποθήκευσης shader.
\subsection{Χώρος υπολογισμού}
Ο χώρος στον οποίο λειτουργεί ένα shader υπολογισμού είναι αφηρημένος. Υπάρχει η έννοια της ομάδας εργασίας. Είναι ο μικρότερος αριθμός από λειτουργίες υπολογισμού τις οποίες μπορεί να εκτελέσει ο χρήστης.
Ο αριθμός των ομάδων εργασίας με τον οποίο μια λειτουργία υπολογισμού εκτελείται, ορίζεται απο τον χρήστη όταν επικαλείται την λειτουργία υπολογισμού. Ο χώρος αυτών των ομάδων είναι τρισδιάστατος, οπότε έχει ένα αριθμό απο ομάδες "Χ","Υ","Ζ". Κάθε ένας απο αυτούς μπορεί να είναι 1, οπότε είναι δυνατή η εκτέλεση λειτουργιών υπολογισμού δυο ή και μίας διάστασης αντί για τρεις διαστάσεις. Αυτό είναι χρήσιμο για την επεξεργασία δεδομένων εικόνας ή γραμμικών πινάκων ενός συστήματος.\\
Η κάθε ομάδα εργασίας μπορεί να αποτελείται απο πολλούς shader υπολογισμού. Αυτό ονομάζεται τοπικό μέγεθος της ομάδας εργασίας. Κάθε shader υπολογισμού έχει τρισδιάστατο τοπικό μέγεθος (το οποίο μπορεί να είναι 1 επιτρέποντας δισδιάστατη ή μονοδιάστατη επεξεργασία). Αυτό ορίζει τον αριθμό των επικλήσεων ενός shader που θα εκτελεστούν σε κάθε ομάδα εργασίας. Για παράδειγμα, αν το τοπικό μέγεθος ενός shader υπολογισμού είναι (128, 1, 1) και εκτελεστεί με έναν αριθμό ομάδων εργασίας (16, 8 , 64) τότε θα έχουμε 1,048,576 ξεχωριστές επικλήσεις shader. Κάθε επίκληση θα έχει ένα σετ από εισόδους που αναγνωρίζουν μοναδικά την κάθε επίκληση. Αυτός ο διαχωρισμός είναι χρήσιμος για διάφορες μορφές συμπίεσης και αποσυμπίεσης εικόνας. Το τοπικό μέγεθος θα είναι το μέγεθος ενός block δεδομένων εικόνας (8x8 για παράδειγμα), ενώ ο αριθμός των ομάδων θα είναι το μέγεθος της εικόνας διαιρούμενο με το μέγεθος του block. Κάθε block κατεργάζεται σαν μια μοναδική ομάδα εργασίας.
\subsection{Υλοποίηση OpenGL}
\subsubsection{Αποστολή}
Ένα αντικείμενο προγράμματος μπορεί να έχει shader υπολογισμού μέσα του. Ο shader υπολογισμού συνδέεται με καταστάσεις shader μέσω κάποιων λειτουργιών απόδοσης. Υπάρχουν δύο λειτουργίες για να ξεκινήσουν οι διαδικασίες υπολογισμού. Χρησιμοποιούν οποιονδήποτε shader υπολογισμού είναι ενεργός. Οι λειτουργίες είναι οι εξής:
\begin{itemize}
\item void glDispatchCompute(GLuint num\_groups\_x, GLuint num\_groups\_u, GLuint num\_groups\_z); - Οι παράμετροι num\_groups\_* ορίζουν τον αριθμό των ομάδων εργασίας, σε τρεις διαστάσεις. Αυτοί οι αριθμοί δεν μπορούν να είναι μηδέν. Υπάρχουν όρια στον αριθμό των ομάδων εργασίας που μπορούν να αποσταλούν.
\item void glDispatchComputeIndirect(GLintptr indirect); - H παράμετρος indirect είναι το αντιστάθμισμα του buffer GL\_DISPATCH\_INDIRECT\_BUFFER. Ισχύουν τα ίδια όρια του αριθμού ομάδων εργασίας, όμως η αποστολή indirect παρακάμπτει τον έλεγχο λαθών του OpenGL. Έτσι, η αποστολή με εκτός ορίων μεγέθους ομάδας εργασίας, μπορεί να προκαλέσει προβλήματα ακόμα και πάγωμα του συστήματος.
\end{itemize}
\subsubsection{Είδοδοι}
Τα shader υπολογισμού  δεν μπορούν να έχουν μεταβλητές καθορισμένες απο τον χρήστη. Τα shader υπολογισμού έχουν τις παρακάτω ενσωματωμένες μεταβλητές εξόδου:
\begin{itemize}
\item in uvec3 gl\_NumWorkGroups; - Αυτή η μεταβλητή περιέχει τον αριθμό των ομάδων εργασίας για την λειτουργία αποστολής
\item in uvec3 gl\_WorkGroupID; - Αυτή η μεταβλητή περιέχει την ισχύουσα ομάδα εργασίας για την επίκληση του shader.
\item in uvec3 gl\_LocalInvocationID; - Αυτή η μεταβλητή περιέχει την ισχύουσα επίκληση του shader μέσα στην ομάδα εργασίας.
\item in uvec3 gl\_GlobalInvocationID; - Αυτή η μεταβλητή αναγνωρίζει μοναδικά την συγκεκριμένη επίκληση του shader υπολογισμού  ανάμεσα σε όλες τις επικλήσεις της κλήσης αποστολής υπολογισμού. Είναι μια συντόμευση για τον μαθηματικό υπολογισμό gl\_WorkGroupID * gl\_WorkGroupSize + gl\_LocalInvocationID;
\item in uint  gl\_LocalInvocationIndex;
\end{itemize}
\subsubsection{Τοπικό μέγεθος}
Το τοπικό μέγεθος ενός shader υπολογισμού ορίζεται απο τον shader, χρησιμοποιώντας μια ειδική δήλωση εισόδου: 
layout(local\_size\_x = X, local\_size\_y = Y, local\_size\_z = Z) in;
Αρχικά, τα τοπικά μεγέθη είναι 1, οπότε αν θέλουμε μονοδιάστατο ή δισδιάστατο χώρο ομάδων εργασίας, μπορούμε να ορίσουμε μόνο το Χ ή το Χ και το Υ. Πρέπει να είναι σταθερές εκφράσεις τιμής μεγαλύτερης του 0. Οι τιμές πρέπει να ορίζονται σε σχέση με τους περιορισμούς που υπάρχουν παρακάτω. Σε αντίθετη περίπτωση προκύπτουν λάθη. Το τοπικό μέγεθος είναι διαθέσιμο στον shader σαν σταθερά, οπότε δεν χρειάζεται να την ορίζουμε εμείς.
\begin{itemize}
\item const uvec3 gl\_WorkGroupSize;
\end{itemize}
\subsubsection{Περιορισμοί}
Ο αριθμός των ομάδων εργασίας που μπορούν να αποσταλούν, ορίζεται από την GL\_MAX\_COMPUTE\_WORK\_GROUP\_COUNT. Αυτή η σταθερά πρέπει να διαβαστεί απο την glGetIntegeri\_v, με τιμές ανάμεσα στο κλειστό όριο [0,2]. Προσπάθεια να καλέσουμε την glDispatchCompute με τιμές που ξεπερνούν το όριο είναι λάθος. Προσπάθεια κλήσης της glDispatchComputeIndirect είναι χειρότερα, μπορεί να διακόψει την λειτουργία του προγράμματος ακόμα και να παγώσει το σύστημα. Σημείωση: ο μικρότερος αριθμός αυτών των τιμών πρέπει να είναι 65535 σε όλους τους άξονες. Αυτό δίνει αρκετό χώρο για εργασία. Υπάρχουν όρια στο τοπικό μέγεθος επίσης. Συγκεκριμένα, υπάρχουν δύο τύποι περιορισμών. 
\begin{itemize}
\item Ο γενικός περιορισμός των διαστάσεων τοπικού μεγέθους, σε συνδυασμό με την GL\_MAX\_COMPUTE\_WORK\_GROUP\_SIZE, όπως και παραπάνω. Η διαφορά είναι οτι ο μικρότερος αριθμός των τιμών είναι πολύ μικρότερος. 1024 για τον Χ και τον Υ, και μόνο 64 για τον Ζ.
\item Ο αριθμός των επικλήσεων μέσα σε μια ομάδα εργασίας. Δηλαδή, το προϊόν των στοιχείων Χ,Υ,Ζ του τοπικού μεγέθους πρέπει να είναι μικρότερο απο GL\_MAX\_WORK\_GROUP\_INVOCATIONS. Η μικρότερη τιμή είναι 1024.
\end{itemize}
Υπάρχει ακόμα ο περιορισμός του ολικού μεγέθους αποθήκευσης για όλες τις κοινές μεταβλητές ενός shader υπολογισμού. Ορίζεται απο την GL\_MAX\_COMPUTE\_SHARED\_MEMORY\_SIZE, που αναφέρεται σε bytes. Η μικρότερη τιμή για το OpenGL είναι 32KB.
\subsection{Υλοποίηση DirectX}
Ένας shader υπολογισμού είναι μια κατάσταση shader υπολογισμού που εξαπλώνει το Microsoft Direct3D 11 πέρα απο τον προγραμματισμό γραφικών. Η τεχνολογία αυτή είναι γνωστή και ως τεχνολογία DirectCompute\cite{computeshaders-4}
\\Όπως όλα τα προγραμματιστικά shader (για παράδειγμα shader γεωμετρίας και κορυφών), ένα shader υπολογισμού είναι σχεδιασμένο να χρησιμοποιεί μια Γλώσσα Υψηλού Προγραμματισμού Shader(HLSL) για το DirectX. Η HLSL, χρησιμοποιείται για το DirectX και μας δίνει την δυνατότητα να δημιουργήσουμε C like shaders για την γραμμή σωλήνων Direct3D. Η HLSL δημιουργήθηκε ξεκινώντας απο το DirectX 9 για την κατασκευή προγραμματιζόμενων τρισδιάστατων γραμμής σωλήνα. Μας δίνει την δυνατότητα να προγραμματίσουμε την γραμμή σωλήνα με τον συνδυασμό οδηγιών assembly, οδηγιών HLSL, και δηλώσεις καθορισμένων λειτουργιών.\\
Ένα shader υπολογισμού προμηθεύει υψηλής ταχύτητας υπολογισμούς γενικού προγραμματισμού, και εκμεταλλεύεται τον μεγάλο αριθμό παράλληλων επεξεργαστών που βρίσκονται στην μονάδα επεξεργασίας γραφικών (GPU). Τα shader υπολογισμού προμηθεύει διαμοιρασμό μνήμης και συγχρονισμό νημάτων, για να επιτρέψει καλύτερες μεθόδους παράλληλου προγραμματισμού. Με την κλήση των μεθόδων ID3D11DeviceContext::Dispatch ή ID3D11DeviceContext::DispatchIndirect γίνεται η εκτέλεση εντολών σε ένα shader υπολογισμού, οι οποίες μπορούν να εκτελεστούν παράλληλα σε πολλά νήματα.
\section{PathScale Enzo}
\chapter{Περισσότερα}

\subsection{Διαφορές CPUs και GPUs}

Ένας απλός τρόπος για να κατανοήσουμε την διαφορά ανάμεσα σε μια CPU και μια GPU είναι να συγκρίνουμε πως αυτές επεξεργάζονται τις διεργασίες. Μια CPU αποτελείται από μερικούς πυρήνες βελτιστοποιημένους για σειριακή επεξεργασία σειράς, ενώ μια GPU έχει μια ογκώδης παράλληλη αρχιτεκτονική που αποτελείται από χιλιάδες μικρότερους, πιο αποδοτικούς πυρήνες σχεδιασμένους για χειρισμό πολλαπλών διαδικασιών ταυτόχρονα.

\begin{figure}[h]
\centering
\includegraphics[scale=1]{cpuvsgpu1}
\caption{Πυρήνες σε CPU και GPU αντίστοιχα}
\end{figure}

Οι GPUs έχουν δημιουργηθεί για πολύ συγκεκριμένες χρήσεις, για παράδειγμα η απεικόνιση γραφικών, πολλαπλασιασμούς πυκνού πλέγματος, απλά φίλτρα επεξεργασίας, κ.α. Είναι πολύ καλές στην διαχείριση μεγάλων καθυστερήσεων γιατί είναι σχεδιασμένες με τρόπο ώστε να ανέχονται καθυστερήσεις υφών, μια διαδικασία 1000+ κύκλων. Οι πυρήνες GPU διαθέτουν πολλά νήματα: όταν ένα νήμα εκτελέσει μια διεργασία υψηλής καθυστέρησης, (για παράδειγμα μια προσπέλαση μνήμης), αυτό το νήμα τίθεται σε αναμονή (ενώ τα υπόλοιπα νήματα συνεχίζουν να εργάζονται), έως ότου η διαδικασία τελειώσει. Αυτό επιτρέπει στις GPUs να διατηρούν τις μονάδες υπολογισμού απασχολημένες περισσότερο από ότι οι παραδοσιακοί πυρήνες.

Οι GPUs δεν αποδίδουν στον χειρισμό διακλαδώσεων, γιατί συνήθως δημιουργούν δεσμίδες από νήματα σε στημόνι, και τα αποστέλλουν στην γραμμή σωλήνα για λόγους οικονομίας εντολών προσπέλασης/αποκωδικοποίησης. Αν τα νήματα συναντήσουν κάποια διακλάδωση, μπορεί να αποκλίνουν, π.χ 2 νήματα σε ένα στημόνι 8 νημάτων μπορεί να ακολουθήσουν την διακλάδωση, ενώ τα άλλα 6 μπορεί να μην την ακολουθήσουν. Τώρα το στημόνι έχει διαιρεθεί σε δύο στημόνια μεγέθους 2 και 6, τα οποία δεν θα εκτελούνται αποδοτικά. Το στημόνι των 2 νημάτων θα εκτελείται με 25\% απόδοση, και το στημόνι των 6 νημάτων θα εκτελείται με 75\% απόδοση. Μπορούμε να φανταστούμε ότι αν μια μονάδα επεξεργασίας συνεχίσει να συναντά επιπλέον διακλαδώσεις, η απόδοση τους συνεχίζει να μειώνεται. Για αυτόν τον λόγο, οι GPUs δεν είναι καλή επιλογή για τον χειρισμό διακλαδώσεων και ο κώδικας που περιέχει αυτές θα πρέπει να εκτελείται στις CPUs.

Γενικότερα, ο κώδικας που δεν θα πρέπει να εκτελείται σε GPUs είναι κώδικας με λίγο παραλληλισμό ή κώδικας με πολλές διακλαδώσεις ή συγχρονισμό, για παράδειγμα βάσεις δεδομένων, λειτουργικά συστήματα, αλγόριθμοι γραφημάτων, κ.α.
\subsubsection{Διαφορές προγραμματιστικού μοντέλου}
Οι πιο σημαντικές διαφορές στο προγραμματιστικό μοντέλο των GPUs είναι ότι δεν υποστηρίζουν διακοπές και εξαιρέσεις. Εκτός από αυτό, δεν υπάρχουν μεγάλες διαφορές μεταξύ της CUDA,OpenCL, και C. Επίσης όπως αναφέραμε, ο κώδικας θα πρέπει να γράφεται με όσο το δυνατόν λιγότερες διακλαδώσεις και συνεχή επικοινωνία μεταξύ των νημάτων.

Πολλά προβλήματα στον πραγματικό κόσμο έχουν πολλές διακλαδώσεις και παρατυπίες. Αλγόριθμοι γραφημάτων, λειτουργικά συστήματα, web browsers, κ.α. Ακόμα και τα γραφικά χρησιμοποιούν όλο και περισσότερο διακλαδώσεις και διαδικασίες γενικού σκοπού, αναγκάζοντας έτσι τις GPUs να γίνονται όλο και περισσότερο προγραμματιζόμενες.


\subsubsection{Ανάγκη χώρου}
Ο πιο σωστός τρόπος για να δούμε την επίδραση του GPGPU είναι να μελετήσουμε τους υπερ-υπολογιστές. Τα μεγάλα ερευνητικά κέντρα (Πανεπιστήμια, Εταιρίες, Κυβερνήσεις) χρειάζονται πολύ μεγάλη υπολογιστική δύναμη για να λύσουν τα πιο δύσκολα υπολογιστικά προβλήματα. Όσο μεγαλύτερο το πρόβλημα, τόσο μεγαλύτερος ο υπολογιστής που χρειάζεται. Παραδείγματα χρήσεων αυτών των υπολογιστών είναι οι διπλώσεις βιομοριακών πρωτεϊνών, η εξομοίωση εκρήξεων και πυρηνικών όπλων, αλλά και οικονομικές συναλλαγές που χρειάζονται για την διακίνηση χρήματος σε όλον τον κόσμο.

Αυτοί οι υπολογιστές συνήθως είναι μεγάλοι, τόσο σε φυσικό μέγεθος όσο και σε υπολογιστικές δυνατότητες. Για παράδειγμα ο TACC Ranger:
\begin{itemize}
\item Αριθμός Κόμβων: 3,936
\item Αριθμός πυρήνων: 62,976
\item Μέγιστη απόδοση: 579.4 TFlops
\item Αριθμός ραφιών: 82
\end{itemize}

Έτσι σε 82 μονάδες έχουμε 580 TFlops υπολογιστικής απόδοσης. Σχετικά με το ίδιο παράδειγμα με χρήση NVIDIA Tesla μονάδων έχουμε
\begin{itemize}
\item 1 Μονάδα Tesla S1070: 4 TFlops
\item 145 μονάδες: 580 TFlops
\item Μονάδες σε 1 ράφι: 42
\item Αριθμός ραφιών: 3.5
\end{itemize}
Δηλαδή για να έχουμε την ίδια απόδοση με την χρήση NVIDIA Tesla μονάδων θα χρειαστούμε το 1/10 του χώρου.

\section{Συμπλέγματα GPU}

\subsection{VirtualCL}
\subsubsection{Επισκόπηση}
Η πλατφόρμα VirtualCL είναι ένα περικάλυμμα για το OpenCL που επιτρέπει τις περισσότερες εφαρμογές να εκμεταλλεύονται με διαφάνεια πολλές OpenCL συσκευές σε ένα σύμπλεγμα, σαν αυτές να βρίσκονταν σε έναν προσωπικό υπολογιστή. Με την πλατφόρμα VirtualCL, οι απομακρυσμένοι κόμβοι εκτελούν διεργασίες για λογαριασμό των προγραμμάτων του ξενιστή.

Το VCL είναι ευέλικτο. Οι εφαρμογές μπορούν να δημιουργήσουν περιεχόμενο OpenCL που αποτελείται από συσκευές από διάφορους κόμβους, ή πολλά περιεχόμενα, κάθε ένα αποτελούμενο από τις συσκευές διαφορετικού κόμβου. Ή οποιονδήποτε συνδυασμό από τα παραπάνω. Άλλες εφαρμογές μπορούν να διαχωριστούν σε διάφορες ανεξάρτητες διεργασίες και νήματα, κάθε μία από τις οποίες εκτελείται σε διαφορετικό σύνολο συσκευών, ενώ ταυτόχρονα χρησιμοποιούν την κοινή μνήμη του συστήματος. 

Πιο εξελιγμένες εφαρμογές μπορούν να επιλέγουν τις συσκευές που θα εκτελούν τις διεργασίες τους, όμως το VCL επιτρέπει μεταβλητές συστήματος στις οποίες ορίζονται πολιτικές για τον διαμοιρασμό των συσκευών. Η αρχική ρύθμιση, είναι κάθε περιεχόμενο που δημιουργείται να περιέχει όλες τις συσκευές ενός κόμβου. Το VCL αποτελείται από τρία μέρη, την βιβλιοθήκη VCL, τον μεσίτη, και τον back-end δαίμονα. 

\begin{figure}[h]
	\includegraphics[width=\linewidth]{gpu-cluster}
	\centering
	\caption{Ένα σύμπλεγμα από υπολογιστικούς κόμβους GPU}
\end{figure}

\subsubsection{Βιβλιοθήκη}
Η βιβλιοθήκη VCL είναι μια βιτρίνα συμπλέγματος για τις εφαρμογές OpenCL. Όταν συνδεθεί με τις εφαρμογές OpenCL, επιτρέπει πρόσβαση με διαφάνεια σε συσκευές OpenCL του συμπλέγματος, κρύβοντας την πραγματική τοποθεσία των συσκευών από τις εφαρμογές που τις καλούν. Η βιβλιοθήκη VCL είναι σχεδιασμένη να εκτελείται με τις περισσότερες εφαρμογές χωρίς επιπλέον παρεμβάσεις, έτσι μπορούμε να επιλέξουμε τον τρόπο που θα χρησιμοποιούνται οι συσκευές στο σύμπλεγμα, μέσω μεταβλητών συστήματος. Η βιβλιοθήκη VCL υποστηρίζει απόλυτα πολυ-νηματικό προγραμματισμό και παρέχει ασφάλεια νημάτων.

Η VCL βιβλιοθήκη χρησιμοποιεί διάφορους αλγόριθμους βελτιστοποίησης. Για παράδειγμα, λόγω της καθυστέρησης δικτύου, η βιβλιοθήκη προσπαθεί να βελτιστοποιήσει την απόδοση της επικοινωνίας διατηρώντας μια ανεξάρτητη βάση δεδομένων από αντικείμενα OpenCL και εκτελεί όσες λειτουργίες είναι δυνατόν στον υπολογιστή-ξενιστή, ώστε να ελαττώσει τον αριθμό των επικοινωνιών στο ελάχιστο.
\subsubsection{Μεσίτης}
Ο μεσίτης είναι μια λειτουργία δαίμονα που εκτελείται σε κάθε υπολογιστή-ξενιστή όπου οι χρήστες μπορούν να εκτελέσουν τις OpenCL εφαρμογές τους. Ο μεσίτης συνδέεται με την βιβλιοθήκη VCL μέσω υποδοχών UNIX. Η ευθύνη του περιέχει:
\begin{enumerate}
\item Παρακολούθηση της λειτουργικότητας και διαθεσιμότητας των συσκευών στο σύμπλεγμα.
\item Αναφορά αυτών των συσκευών στις εφαρμογές που τις ζητάνε.
\item Έξυπνη κατανομή των συσκευών για τις εφαρμογές OpenCL όταν γίνεται δημιουργία περιεχομένου, για παράδειγμα η προσπάθεια να ταιριάξει ο αριθμός των συσκευών με τον αριθμό τον κόμβων που έχουν σύνολο τον αριθμό που ζητήθηκε απο την εφαρμογή.
\item Πιστοποίηση αυθεντικότητας, δρομολόγηση, και εξασφάλιση της ποιότητας των μηνυμάτων μεταξύ των εφαρμογών. 
\end{enumerate}

\subsubsection{Back-end Δαίμονας}
Η υπηρεσία του back-end δαίμονα εκτελείται σε κάθε κόμβο συμπλέγματος όπου υπάρχουν συσκευές OpenCL και υποστηρίζονται από κατάλληλους οδηγούς. Ο δαίμονας χρησιμοποιεί οποιαδήποτε OpenCL βιβλιοθήκη είναι διαθέσιμη στους κόμβους για να εκτελέσει πυρήνες για λογαριασμό των εφαρμογών πελάτη. Εξομοιώνει όλες τις διαδικασίες OpenCL όπως αυτές ζητούνται από την βιβλιοθήκη VirtualCL. Για λόγους ασφάλειας, όσο οι συσκευές GPU και οι οδηγοί δεν επιτρέπουν διαφανής προτίμηση, οι συσκευές OpenCL δεν κοινοποιούνται από το back-end στις υπόλοιπες εφαρμογές. Κάθε συσκευή κατανέμεται σε μόνο μια εφαρμογή κάθε στιγμή.


\subsubsection{Προβλήματα}
Όταν εκτελούμε πυρήνες OpenCL σε απομακρυσμένες συσκευές, η καθυστέρηση δικτύου είναι ο κύριος παράγοντας περιορισμού. Η ελαχιστοποίηση του αριθμού των "ταξιδιών" για βασικές λειτουργίες OpenCL είναι το πρώτο βήμα για να αντιμετωπίσουμε αυτό το πρόβλημα. Γιαυτό αναπτύχθηκε το SuperCL, όπου μια σειρά απο πυρήνες και λειτουργίες μνήμης αποστέλλονται στις συσκευές ενός κόμβου συμπλέγματος, συνήθως μόνο σε ένα "ταξίδι". Όταν είναι απαραίτητο, επιτρέπεται η επικοινωνία με τον ξενιστή, αλλά με ασύγχρονο τρόπο, για να αποφευχθεί περαιτέρω καθυστέρηση. Το εύρος ζώνης συχνότητας είναι επίσης περιοριστικός παράγοντας όταν περιέχονται τεράστιες διεργασίες I/O, και αυτό αντιμετωπίζεται από το SuperCL επιτρέποντας την αρχικοποίηση των buffer από αρχεία back-end, και την αποθήκευση των αποτελεσμάτων σε αρχεία back-end.

\
\chapter{Εφαρμογές}
\section{Εισαγωγή}
Τα τελευταία περίπου 20 χρόνια οι εταιρίες παραγωγής υλικού γραφικών έχουν εστιάσει στην προσπάθεια να παράγουν γρήγορες μονάδες γραφικής επεξεργασίας (GPU), ειδικότερα για την κοινότητα των gamer. Αυτό έχει ως αποτέλεσμα πρόσφατα να δημιουργηθούν συσκευές οι επιδόσεις των οποίων ξεπερνούν τις κεντρικές μονάδες επεξεργασίας (CPU), σε συγκεκριμένες εφαρμογές, ειδικότερα σε μετρήσεις εκατομμυρίων εντολών το δευτερόλεπτο (MIPS). Έτσι, καθιερώθηκε μια κοινότητα για να αξιοποιήσει αυτήν την μεγάλη δύναμη των GPU για υπολογισμούς γενικής χρήσης(GPGPU). Τα τελευταία δύο χρόνια έχουν εξαλειφθεί οι περισσότεροι περιορισμοί που υπήρχαν όσον αφορά το σετ εντολών και την διαχείριση μνήμης, με την ενσωμάτωση ενοποιημένων υπολογιστικών μονάδων στις κάρτες γραφικών, δίνοντας έτσι την δυνατότητα στους προγραμματιστές να δημιουργήσουν ένα πλήθος από προγράμματα με εφαρμογές σε πολλούς τομείς.
\section{Κρυπτογράφηση}
Στο πεδίο της ασύμμετρης κρυπτογράφησης, η ασφάλεια όλων των πρακτικών κρυπτοσυστημάτων βασίζεται στην δυσκολία υπολογισμού προβλημάτων, εξαρτημένη από την επιλογή των παραμέτρων. Με την όποια αύξηση των παραμέτρων όμως (συνήθως στο εύρος 1024-4096 bits), οι υπολογισμοί γίνονται όλο και πιο απαιτητικοί για τον εκάστοτε επεξεργαστή. Σε σύγχρονο υλικό, ο υπολογισμός μιας μονής εντολής κρυπτογράφησης δεν είναι κρίσιμος, όμως σε ένα σύστημα επικοινωνίας πολλών-προς-ένα, για παράδειγμα ένας κεντρικός server στο κέντρο δεδομένων μιας εταιρίας, μπορεί να αντιμετωπίσει ταυτόχρονα εκατοντάδες η και χιλιάδες ταυτόχρονες συνδέσεις και εντολές κρυπτογράφησης.

Ως αποτέλεσμα, η πιο συνήθης λύση για ένα τέτοιο σενάριο είναι η χρήση καρτών επιτάχυνσης κρυπτογράφησης. Λόγω της μικρής αγοράς, η τιμή τους φτάνει συνήθως αρκετά χιλιάδες ευρώ η δολάρια.
Τελευταία, η ερευνητική κοινότητα έχει αρχίσει να εξερευνά τεχνικές για επιτάχυνση των αλγορίθμων κρυπτογράφησης με χρήση της GPU. \cite{cryptography-1}

\begin{figure}[h]
\centering
\includegraphics[scale=0.50]{crypto-1}
\caption{Μηχανή κρυπτογράφησης German Lorenz, χρησιμοποιήθηκε στον δεύτερο παγκόσμιο πόλεμο για να κρυπτογραφεί μηνύματα για προσωπικό πολύ υψηλής σημασίας\cite{figure-7}}
\end{figure}

\subsection{Κρυπτογράφηση συμμετρικού κλειδιού}
Η κρυπτογραφία συμμετρικού κλειδιού αναφέρεται στις κρυπτογραφικές μεθόδους στις οποίες ο αποστολέας και ο αποδέκτης μοιράζονται το ίδιο κλειδί (ή σε πιο σπάνια περίπτωση, όταν τα κλειδιά είναι διαφορετικά, αλλά βρίσκονται εύκολα μέσω απλού υπολογισμού). Αυτή ήταν η μόνη μορφή κρυπτογραφίας γνωστή μέχρι το 1976.
Οι κρυπτογράφοι συμμετρικών κλειδιών υλοποιούνται σαν block ή κρυπτογράφοι ροής. Ένας κρυπτογράφος block κρυπτογραφεί τα δεδομένα σε blocks αντί για τον κάθε χαρακτήρα ξεχωριστά, δηλαδή με τον τρόπο που λειτουργεί ο κρυπτογράφος ροής.\cite{cryptography-2}


\begin{figure}[h]
\centering
\includegraphics[scale=0.50]{crypto-2}
\caption{Γαλλική μηχανή κρυπτογράφησης σε σχήμα βιβλίου του 16ου-αιώνα\cite{figure-8}}
\end{figure}

Το πρότυπο κρυπτογράφησης δεδομένων (DES) και το ανεπτυγμένο πρότυπο κρυπτογράφησης (AES) είναι κρυπτογράφοι block σχεδιασμένοι από την Αμερικάνικη κυβέρνηση. Αν και έχει σταματήσει να υποστηρίζεται σαν επίσημο πρότυπο, η χρήση του DES είναι ακόμα αρκετά δημοφιλής. Χρησιμοποιείται για μια μεγάλη σειρά από εφαρμογές, όπως κρυπτογράφηση ATM, ασφαλής απομακρυσμένη πρόσβαση, ασφαλής επικοινωνία ηλεκτρονικού ταχυδρομείου, κ.α.

Οι κρυπτογράφοι ροής, σε αντίθεση με τους block, δημιουργούν μια μεγάλη σειρά από υλικά κλειδιών, που συνδυάζονται με το απλό κείμενο χαρακτήρα προς χαρακτήρα. Σε έναν κρυπτογράφο ροής, το αποτέλεσμα παράγεται βασισμένο σε μια κρυμμένη κατάσταση που αλλάζει όσο λειτουργεί ο κρυπτογράφος. Αυτή η εσωτερική κατάσταση αρχικοποιείται με την χρήση ενός κρυμμένου κλειδιού. Παράδειγμα γνωστού κρυπτογράφου ροής είναι ο RC4.\cite{cryptography-3}

Οι κρυπτογραφικές διαδικασίες hash είναι ένας τρίτος τύπος αλγόριθμου κρυπτογράφησης. Δέχονται σαν είσοδο οποιαδήποτε μήκος κειμένου, και δίνουν αποτέλεσμα ένα μικρό κείμενο συγκεκριμένου μήκους που μπορεί να χρησιμοποιηθεί σαν ψηφιακή ταυτότητα. Όσον αφορά τις καλές διαδικασίες hash, ένας εισβολέας δεν μπορεί να βρει δύο μηνύματα που παράγουν το ίδιο hash. \cite{cryptography-4}To MD5 για παράδειγμα, ενώ χρησιμοποιείται εδώ και πολλά χρόνια, έχει αποδειχτεί ότι είναι εύκολο να παραβιαστεί. Το SHA-1 είναι περισσότερο διαδεδομένο ως ασφαλής διαδικασία hash κρυπτογράφησης, αλλά αναλυτές έχουν ανακαλύψει επιθέσεις για την παραβίαση του. Το 2012, μια καινούρια διαδικασία κρυπτογράφησης hash ανακαλύφθηκε, η SHA-3, η οποία χρησιμοποιεί τον αλγόριθμο Keccak για την υλοποίηση της. Σε παρακάτω υπο-ενότητα θα μελετήσουμε την χρήση προγράμματος για παραβίαση των διαδικασιών κρυπτογράφησης hash.\cite{cryptography-5}
\subsection{Κρυπτογράφηση δημοσίου κλειδιού}
Τα κρυπτοσυστήματα συμμετρικών κλειδιών χρησιμοποιούν το ίδιο κλειδί για την κρυπτογράφηση και την αποκρυπτογράφηση ενός μηνύματος, αν και ένα μήνυμα ή ομάδα μηνυμάτων μπορεί να έχουν διαφορετικό κλειδί από τα άλλα. Ένα σημαντικό μειονέκτημα των συμμετρικών κρυπτογραφημάτων είναι η διαχείριση κλειδιών, ώστε αυτά να χρησιμοποιηθούν με ασφάλεια. 

\begin{figure}[h]
\centering
\includegraphics[width=\linewidth]{crypto-3}
\caption{Η μηχανή κρυπτογράφησης Enigma, χρησιμοποιήθηκε από τον Γερμανικό στρατό και τις πολιτικές αρχές από τα τέλη του 1920 μέχρι και τον δεύτερο παγκόσμιο πόλεμο, παρείχε ένα πολύπλοκο ηλεκτρο-μηχανικό πολύ-αλφαβητικό κρυπτογράφημα. H αποκρυπτογράφηση του αλγόριθμου αποδείχτηκε μεγάλης σημασίας για την νίκη των συμμάχων.\cite{figure-9}}
\end{figure}

\subsection{Hashcat}
Για λόγους έρευνας της εργασίας επιλέχτηκε το πρόγραμμα oclHashcat ώστε να μελετήσουμε την απόδοση των κρυπτογραφικών προγραμμάτων με χρήση GPGPU. Το Hashcat είναι το πιο γρήγορο πρόγραμμα επαναφοράς κωδικών. Είναι ελεύθερο, αν και κλειστού κώδικα. Υπάρχουν εκδόσεις για Linux,OSX,και Windows, και υλοποιήσεις για CPU ή GPU. Το hashcat υποστηρίζει μεγάλο αριθμό από hashing αλγορίθμους, συμπεριλαμβανομένου MD4,MD5,SHA, Unix Crypt, κ.α
\subsubsection{Υλοποιήσεις}
Το Hashcat διανέμεται σε δύο εκδόσεις
\begin{itemize}
\item Hashcat - Ένα εργαλείο επαναφοράς κωδικών για CPU
\item oclHashcat - Επαναφορά κωδικών με χρήση GPGPU
\end{itemize}

Πολλοί από τους αλγόριθμους που υποστηρίζει το hashcat μπορούν να βρεθούν σε μικρότερο χρόνο με την χρήση της GPU-επιτάχυνσης του oclHashcat(όπως MD5,SHA1, κ.α) Όμως, δεν επιταχύνονται όλοι οι αλγόριθμοι απο την χρήση GPU. Το Bcrypt είναι ένα τέτοιο παράδειγμα. Λόγω παραγόντων όπως διακλαδώσεις δεδομένων, σειριακοποίηση, μνήμη, κ.α, το oclHashcat δεν είναι απόλυτος αντικαταστάτης του Hashcat.

\subsubsection{Επιθέσεις}
Το Hashcat υποστηρίζει πολλούς τύπους επιθέσεων για την επαναφορά πολύπλοκων κωδικών από τον χώρο κλειδιού ενός hash. Αυτοί οι τύποι περιλαμβάνουν:
\begin{itemize}
\item Επίθεση Brute-Force
\item Επίθεση Combinator
\item Επίθεση λεξικού
\item Επίθεση αποτυπώματος
\item Υβριδική επίθεση
\item Επίθεση μάσκας
\item Επίθεση Permutation
\item Επίθεση βασισμένη σε κανόνες
\item Επίθεση με χρήση πίνακα
\item Επίθεση Toggle-Case
\end{itemize}


\subsection{Κρυπτο-νομίσματα}
Τα κρυπτονομίσματα είναι μια μορφή συναλλαγής με χρήση κρυπτογραφίας για ασφάλεια και τον έλεγχο της δημιουργίας καινούριων νομισμάτων. Το πρώτο κρυπτο-νόμισμα που δημιουργήθηκε ήταν το Bitcoin το 2009. Από τότε, πολλά κρυπτονομίσματα έχουν δημιουργηθεί. Ένα χαρακτηριστικό τους είναι ότι δεν υπάρχει κεντρικός έλεγχος, σε αντίθεση με άλλα συστήματα ηλεκτρονικού χρήματος όπως το Paypal. Ακόμα ένα χαρακτηριστικό είναι ότι οι συναλλαγές καταγράφονται δημόσια, για παράδειγμα στο Bitcoin, οι συναλλαγές καταγράφονται στην block αλυσίδα. \cite{cryptography-6}

\begin{figure}[h]
\centering
\includegraphics[width=0.3\linewidth]{bitcoin}
\caption{Λογότυπο Bitcoin, το πρώτο και το πιο επιτυχημένο κρυπτο-νόμισμα\cite{figure-10}}
\end{figure}

Η δημιουργία κρυπτονομισμάτων γίνεται μέσω ειδικών προγραμμάτων που υπολογίζουν συγκεκριμένα hashes. Τα κρυπτονομίσματα χρησιμοποιούν έναν αλγόριθμο "Proof of work" ώστε να γνωρίζουν ότι ο χρήστης έχει εξορύξει τα νομίσματα που αυτός αναφέρει. (π.χ Scrypt, SHA-256, κ.α). Η εξόρυξη ενός block είναι μια διαδικασία που μοιάζει πολύ με την επαναφορά κωδικών. Για αυτό οι GPUs, είναι πιο αποδοτικές απο τις CPUs - όταν μια τυπική CPU έχει μέχρι 8 πυρήνες, μια GPU μπορεί να έχει εκαντοντάδες, όπου η καθεμία απο αυτές υπολογίζει ένα διαφορετικό "hash".

Υπάρχουν δύο τρόποι για εξόρυξη εικονικών νομισμάτων.
\begin{itemize}
\item Ατομική εξόρυξη - Η ατομική προσπάθεια για εξόρυξη νομισμάτων με χρήση μιας η περισσότερων GPU
\item Εξόρυξη κοινοπραξίας - Η προσπάθεια από πολλούς ανθρώπους να εξορύξουν το ίδιο block, διαιρώντας τα κέρδη.
\end{itemize}
Οι ταχύτητες της εξόρυξης μετρώνται σε KH/s (kilohashes ανα δευτερόλεπτο) και MH/s (megahashes ανα δευτερόλεπτο). Για παράδειγμα, μια Radeon 4870 εξορύσσει ένα νόμισμα SHA-256 με ταχύτητα 80-100 MH/s, ενώ ένα νόμισμα Scrypt με ταχύτητα μόλις 130-140 KH/s.

H χρήση των GPUs για εξόρυξη εικονικών νομισμάτων, οδήγησε την αγορά των καρτών γραφικών σε αύξηση των τιμών, αφού για πολύ καιρό ήταν μια αποδοτική λύση σε αντίθεση με την χρήση CPUs. Πρόσφατα όμως, με την παρουσίαση εξειδικευμένων καρτών (ASICS - Application Specific Integrated Circuit), η χρήση των GPUs για εξόρυξη κρυπτο-νομισμάτων άρχισε να μειώνεται, αφού η απόδοση τους είναι έως και 1000 φορές μικρότερη. Αυτό ήταν επίσης ένας λόγος για τον απλό χρήστη να απομακρυνθεί από την εξόρυξη νομισμάτων, αφού δημιουργήθηκαν μικρότερες κοινοπραξίες που απαιτούν μεγαλύτερη υπολογιστική δύναμη.
\section{Βιοπληροφορική}
\subsection{Εισαγωγή}
Η συνεχής αύξηση της ποσότητας βιολογικών δεδομένων, η ανάγκη για ανάλυση τους και το συνεχές ενδιαφέρον από την επιστημονική κοινότητα για την κατανόηση των δομικών λειτουργιών των βιολογικών μορίων, αποτέλεσαν τους κύριους λόγους για την ανάπτυξη της βιοπληροφορικής. Για να κατανοήσουμε τις κυτταρικές και βιομοριακές λειτουργίες, τα βιολογικά δεδομένα πρέπει να συνενωθούν για να σχηματίσουν μια ακριβής εικόνα. Οι ερευνητές της βιοπληροφορικής, έχουν αναπτύξει υπολογιστικές τεχνικές για την επεξεργασία των βιολογικών δεδομένων, όπως νουκλεοτιδικές αλληλουχιών, αλληλουχίες αμινο οξέων, τρισδιάστατων δομών, όπως επίσης βιολογικών σημάτων και εικόνων. Μεγάλες ερευνητικές προσπάθειες του πεδίου συμπεριλαμβάνουν αναγνώριση προτύπων, ευθυγράμμιση αλληλουχιών, ανάλυση πρωτεϊνικών δομών, φυλογενητική ανάλυση, μοριακή δυναμική, ανάλυση γονιδιώματος, σχεδιασμός φαρμάκων και ανάπτυξη φαρμάκων. Επίσης, υπάρχουν προφητικές τεχνικές ειδικές για τις εκφράσεις γονιδίων, και την αλληλεπίδραση πρωτεϊνών.\\
\begin{figure}[h]
\centering
\includegraphics[scale=0.25]{bioinformatics}
\caption{Βιολογία και πληροφορική}
\end{figure}
Η Βιοπληροφορική παίζει μεγάλο ρόλο σε πολλές πτυχές της βιολογίας. Στην πειραματική μοριακή βιολογία, οι τεχνικές βιοπληροφορικής όπως επεξεργασία εικόνας και σήματος, επιτρέπει την εξόρυξη χρήσιμων αποτελεσμάτων από μεγάλο όγκο δεδομένων. Στο πεδίο της γενετικής και γονιδιωματικής, συμβάλλει στην αλληλουχία και υποσημείωση γονιδιωμάτων και την παρατήρηση των μεταλλάξεων τους.\cite{bioinformatics-1} Παίζει μεγάλο ρόλο στην εξόρυξη τεχνικών όρων και στην κατασκευή βιολογικών και γονιδιακών οντολογιών για την οργάνωση και αναζήτηση βιολογικών δεδομένων. Έχει επίσης μεγάλο ρόλο στην ανάλυση των γονιδίων και στην ρύθμιση πρωτεϊνών. Τα εργαλεία της Βιοπληροφορικής συμβάλουν στην σύγκριση γενετικών και γονιδιακών δεδομένων και γενικότερα στην κατανόηση των αναπτυξιακών πτυχών της μοριακής βιολογίας. Σε πιο εσωτερικό επίπεδο, συμβάλει στην ανάλυση και κατηγοριοποίηση των βιολογικών διαδρόμων και δικτύων τα οποία είναι σημαντικό κομμάτι της συστεμικής βιολογίας. Στην Δομική βιολογία, συμβάλει στην εξομοίωση και μοντελισμό του DNA, RNA, και δομές πρωτεϊνών όπως και μοριακών αλληλεπιδράσεων.\\

\subsection{Μοριακή δυναμική}
Η Βιοπληροφορική είναι ένα επιστημονικό πεδίο που εστιάζει στην εφαρμογή της τεχνολογίας υπολογιστών στην διαχείριση βιολογικών δεδομένων. Με το πέρασμα του χρόνου, οι εφαρμογές βιοπληροφορικής έχουν χρησιμοποιηθεί για να αποθηκεύσουν, αναλύσουν και να ενσωματώσουν βιολογικές και γενετικές πληροφορίες, χρησιμοποιώντας ένα μεγάλο εύρος μεθοδολογιών. Μια από τις πλέον γνωστές τεχνικές για την κατανόηση των φυσικών κινήσεων των ατόμων και των μορίων, είναι η μοριακή δυναμική. Η μοριακή δυναμική είναι μια μέθοδος εξομοίωσης των φυσικών κινήσεων των ατόμων και των μορίων κάτω από συγκεκριμένες συνθήκες. Έχει ρόλο κλειδί σε επιστήμες όπως η βιολογία, η χημεία, η φυσική, ιατρική. Λόγω της πολυπλοκότητας τους, οι υπολογισμοί της μοριακής δυναμικής χρειάζονται μεγάλες ποσότητες μνήμης και υπολογιστικής δύναμης, και για αυτό η εκτέλεση τους είναι συχνά μεγάλο πρόβλημα.\cite{bioinformatics-3} \\
Οι εξομοιώσεις της μοριακής δυναμικής χρησιμοποιούν πολύπλοκους αριθμητικούς υπολογισμούς, που πολλές φορές οδηγούν σε αριθμητικά λάθη. Πριν την ανακάλυψη του προγραμματισμού γενικής χρήσης, οι GPU χρησιμοποιούνταν μόνο για διαδικασίες απεικόνισης των μοριακών δομών, και η εκτέλεση των αλγορίθμων μοριακής δυναμικής μπορούσε να διαρκέσει από ώρες, έως και μέρες. Η λύση προήλθε από το GPGPU, καθώς οι GPU έχουν πολλές αριθμητικές μονάδες που μπορούν να εκτελεστούν παράλληλα. \\
Στο πεδίο της μοριακής δυναμικής, έχουν αναπτυχθεί πολλές εφαρμογές εφαρμογές βασισμένα στο GPGPU, που υποστηρίζουν εξομοιώσεις σε πολλαπλές μονάδες. Αυτή η καινοτομία δημιουργεί ευκαιρίες για το μέλλον, ειδικά για μικρότερες ερευνητικές ομάδες. Μειώνει τον χρόνο που απαιτείται για διαδικασίες και τα απαραίτητα κονδύλια για έρευνα-ανάπτυξη, προάγει την ανάπτυξη καινούριων εφαρμογών και την επιστημονική πρόοδο.\cite{bioinformatics-4}
\subsubsection{Μετάβαση από CPU σε GPU}
Η διαφορά στην αρχιτεκτονική μεταξύ CPU και GPU, είναι ότι στην τελευταία είναι δυνατή η εκτέλεση πολλαπλών παράλληλων διεργασιών, κάτι που επιτρέπει την καλύτερη εκτέλεση πολύπλοκων αλγορίθμων και καλύτερη διαχείριση μεγάλου όγκου δεδομένων. Επίσης, μια μονάδα επεξεργασίας γραφικών έχει λιγότερες ενεργειακές απαιτήσεις, έτσι η δημιουργία υπερ-υπολογιστών με χρήση GPU εξαλείφει την ανάγκη για τεράστιους χώρους γεμάτους με υπολογιστές. Η εγκατάσταση μιας επιπλέον μονάδας, αντιγράφει τον παραλληλισμό του προγραμματισμού, χωρίς καμία επιπλέον ενέργεια. Επιπλέον, η GPU έχει εντυπωσιακές δυνατότητες υπολογισμού floating point και μεγάλο εύρος ζώνης μνήμης, δίνοντας την δυνατότητα για βελτιστοποιημένη πρόσβαση στην μνήμη, ελεγχόμενη εκτέλεση επιλογών, και διαχείριση πόρων, με χρήση λίγων γραμμών κώδικα. Συγκεκριμένα, οι εφαρμογές μοριακής δυναμικής, κβαντικής χημείας, η απεικόνιση των αποτελεσμάτων τους, τρέχουν μέχρι και 5 φορές πιο γρήγορα. \\
Από την αρχή του GPU προγραμματισμού μέχρι και σήμερα, η προγραμματιστική ανάπτυξη συνεχίζεται αδιάκοπα. Ο αριθμός των εφαρμογών βασισμένων σε αρχιτεκτονικές GPU, φτάνουν τις 200, το οποίο είναι αύξηση της τάξεως πάνω από 60\% μέσα σε δύο χρόνια. Οι καλύτερες εφαρμογές βασισμένες σε GPGPU έχουν σχεδιαστεί για την μοριακή δυναμική, τον σχεδιασμό φαρμάκων, κβαντική χημεία, το κλίμα, την φυσική σύμφωνα με την Nvidia\cite{bioinformatics-2}
\subsection{GPUGRID.net}
\epigraph{"I hope mankind will acknowledge people like you, its real heroes."}{Grzegorz Granowski, Volunteer \& Donor}
Το GPUGRID είναι ένα εθελοντικό κατανεμημένο σύστημα, το οποίο στοχεύει στην βιοϊατρική έρευνα από το πανεπιστήμιο Universitat Pompeu Fabra της Ισπανίας. Το GPUGRID αποτελείται από πολλές μονάδες επεξεργασίας γραφικών, που συνεργάζονται μεταξύ τους για να παραδώσουν υψηλών επιδόσεων εξομοιώσεις βιομορίων. Οι μοριακές εξομοιώσεις πού εκτελούνται από τους εθελοντές του, αποτελούν μερικούς απο τους πιο συνήθης τύπους εξομοιώσεων που εκτελούνται απο τους επιστήμονες του πεδίου, αλλά ταυτόχρονα είναι από τους πιο απαιτητικούς σε υπολογιστική δύναμη και συνήθως απαιτούν υπερ-υπολογιστές.\\
\begin{figure}[h]
\centering
\includegraphics[scale=0.75]{gpugrid}
\caption{Βιολογία και πληροφορική}
\end{figure}\\
Το σύστημα ερευνά μεταξύ άλλων τα παρακάτω προβλήματα
\begin{itemize}
\item Εξομοίωση της ωρίμανσης πρωτεολυτικών του HIV - Μια απο τις πιο σημαντικές πτυχές της ωρίμανσης του HIV είναι το πώς η πρωτεΐνη "ψαλιδιών", δημιουργείται. Η απάντηση σε αυτό το ερώτημα χρειάζεται εξομοιώσεις μοριακής δυναμικής στο όριο των μοντέρνων υπολογιστικών δυνατοτήτων. Το GPUGRID μας επιτρέπει να λύσουμε αυτο το πρόβλημα και έχουμε καταφέρει να δείξουμε οτι τα πρώτα "ψαλίδια" κόβονται απο το "σκοινί" που είναι δεμένα. Αυτό το γεγονός συμβαίνει στην αρχή της ωρίμανσης, και αν σταματήσουμε την ωρίμανση των πρωτεολυτικών, τότε θα σταματήσουμε και την ωρίμανση του HIV σαν σύνολο.\cite{gpugrid-1}
\item Ανακάλυψη του ρόλου των μεμβρανών λιπιδίων στην δραστηριότητα ενζύμων.
\item Μοριακή εξομοίωση αισθητήρων ντοπαμίνης κάτω απο φυσιολογικές ιονικές δυνάμεις.
\item Αποκάλυψη των μηχανισμών αντίδρασης φαρμάκων καρκίνου παχέος εντέρου - Ο καρκίνος είναι βασικά ή ανεξέλεγκτη ανάπτυξη ιστών και εισβολή από μεταλλαγμένα κύτταρα σε έναν οργανισμό. Σε αντίθεση με τις παραδοσιακές χημειοθεραπείες ή ραδιοθεραπείες, οι νεότερες θεραπείες στοχεύουν σε συγκεκριμένους στόχους κακοήθων κυττάρων. Αυτό επιτυγχάνεται με τον εντοπισμό ορισμένων πρωτεϊνών που εκφράζονται διαφορικά σε ογκογεννητικά κύτταρα. Με την βοήθεια του GPUGRID, επιτυγχάνεται η επεξήγηση των μοριακών μηχανισμών που συμβαίνουν στα μεταλλαγμένα μόρια των κυττάρων.
\end{itemize}
Η εκτέλεση του GPUGRID στις GPUs, καινοτομεί στον εθελοντικό υπολογισμό, παραδίδοντας εφαρμογές υπερ-υπολογιστών, σε υποδομές χαμηλού κόστους. Η απόδοση των μονάδων γραφικής επεξεργασίας, καταγράφεται και συγκρίνεται σε σχέση με άλλους χρήστες, ανάλογα με την διάρκεια ολοκλήρωσης των WU (Work Units)\cite{gpugrid-2}.\\
\begin{figure}[h]
\centering
\includegraphics[scale=0.75]{gpugrid-charts}
\caption{Βιολογία και πληροφορική}
\end{figure}

\subsection{Προγράμματα}
\begin{apptable}{Βιοπληροφορική}{bioinformatics}
BarraCUDA & Λογισμικό χαρτογράφησης ακολουθίας & 6-10x & ΝΑΙ \\ \hline
CUDASW++ & Λογισμικό ανοιχτού κώδικα για αναζητήσεις Smith-Waterman σε πρωτεϊνικές βάσεις δεδομένων με χρήση GPUs & 10-50x & ΝΑΙ \\ \hline
CUSHAW & Ευθυγραμμιστής παράλληλων μικρών προσπελάσεων  & 10x & ΝΑΙ \\ \hline
G-BLASTN & Επιταχυνόμενο από GPU εργαλείο ευθυγράμμισης νουκλεοτιδίων βασισμένο στο ευρέως διαδεδομένο NCBI-BLAST & 4-15x & \\ \hline
GPU-BLAST & Τοπική αναζήτηση με γρήγορους ευρετικούς  αλγόριθμους k-tuple & 3-4x & \\ \hline
mCUDA-MEME & Πολύ γρήγορη κλιμακωτή ανακάλυψη μοτίβων βασισμένη στο MEME & 4-10x & ΝΑΙ \\ \hline
MUMmer GPU & Πρόγραμμα υψηλής απόδοσης τοπικής ευθυγράμμισης ακολουθίας & 3-10x & \\ \hline
NVBIO & Βιβλιοθήκη ανοιχτού κώδικα C++ αποτελούμενη από στοιχεία επαναχρησιμοποιήσιμα σχεδιασμένα για να επιταχύνουν εφαρμογές βιοπληροφορικής με χρήση CUDA. & 4-5x & ΝΑΙ \\ \hline
NVBowtie & Μια σε μεγάλο βαθμό ολοκληρωμένη εφαρμογή του ευθυγραμμιστή Bowtie2 πάνω από το NVBIO & 2.75x-8.35x & ΝΑΙ \\ \hline
PEANUT & Καθορισμός ανάγνωσης για ακολουθίες DNA ή RNA σε γνωστή αναφορά γονιδιώματος. & 10x & \\ \hline
REACTA & Ο σκοπός του REACTA είναι ο ποσοτικός προσδιορισμός της συμβολής της γενετικής διακύμανσης στην φαινοτυπική διακύμανση σε πολύπλοκα χαρακτηριστικά. & 2-4x & ΝΑΙ \\ \hline
SeqNFind & Ακολουθία επόμενης γενιάς και συγκρίσεις γονιδιωμάτων & 400x & ΝΑΙ \\ \hline
SOAP3 & Λογισμικό βασισμένο σε GPU για ευθυγράμμιση μικρών προσπελάσεων με αναφορά ακολουθίας. Μπορεί να βρει όλες τις ευθυγραμμίσεις με k ασυμφωνίες, όπου το k είναι ένας αριθμός απο το 0 έως το 3 & 10x & ΝΑΙ \\ \hline
SOAP3-dp & Πολύ γρήγορο εργαλείο βασισμένο σε GPU για ευθυγραμμίσεις μικρών προσπελάσεων μέσω δυναμικού προγραμματισμού υποβοηθούμενου από ευρετήριο & 28-64x & ΝΑΙ \\ \hline
UGENE & Λογισμικό ανοιχτού κώδικα Smith-Waterman για SSE/CUDA, επαναλήψεις βασισμένες σε πίνακα δεικτών & 6-8x & ΝΑΙ \\ \hline
WideLM & Ταιριάζει πολυάριθμα γραμμικά μοντέλα σε μια σταθερή σχεδίαση και απάντηση & 150x & ΝΑΙ \\ \hline
\end{apptable}


%\begin{apptable}{Μοριακή δυναμική}{molecular}
%ACEMD & GPU simulation of molecular mechanics force fields, implicit and explicit solvent & & \\ \hline
%AMBER & Suite of programs to simulate molecular dynamics on biomolecule & & \\ \hline
%CHARMM & MD package to simulate molecular dynamics on biomolecule & & \\ \hline
%DESMOND & High-speed molecular dynamics simulations of biological systems & & \\ \hline
%DL-POLY & Simulate macromolecules, polymers, ionic systems, etc on a distributed memory parallel computer & & \\ \hline
%ESPResSo & Highly versatile software package for performing and analyzing scientific Molecular Dynamics. & & \\ \hline
%Folding@Home & A distributed computing project that studies protein folding, misfolding, aggregation, and related diseases & & \\ \hline
%GPUGrid.net & A distributed computing project that uses GPUs for molecular simulations & & \\ \hline
%GROMACS & Simulation of biochemical molecules with complicated bond interactions & & \\ \hline
%HALMD & Large-scale simulations of simple and complex liquids & & \\ \hline
%HOOMD-Blue & Particle dynamics package written grounds up for GPUs & & \\ \hline
%LAMMPS & Classical molecular dynamics package & & \\ \hline
%NAMD & Designed for high-performance simulation of large molecular systems & & \\ \hline
%OpenMM & Library and application for molecular dynamics for HPC with GPUs & & \\ \hline
%\end{apptable}
\section{Κβαντική χημεία}
\subsection{Εισαγωγή}
Η κβαντική χημεία είναι μια διακλάδωση της χημείας όπου η κύρια εστίαση είναι η εφαρμογή κβαντικής μηχανικής σε μοντέλα φυσικής και σε πειράματα  χημικών συστημάτων.
Η ιστορία της κβαντικής χημείας ουσιαστικά ξεκινάει με την ανακάλυψη των καθοδικών ακτίνων από τον Michael Faraday, την υπόδειξη του 1877 απο τον Ludwig Boltzmann ότι τα επίπεδα ενέργειας σε ένα σύστημα φυσικής μπορεί να είναι διακριτά, και την κβαντική υπόθεση του 1900 από τον Max Planck ότι κάθε ατομικό σύστημα ακτινοβολίας ενέργειας μπορεί να διαιρεθεί σε έναν αριθμό διακριτών στοιχείων έτσι ώστε κάθε από αυτά τα στοιχεία ενέργειας να είναι αναλογικό στην συχνότητα με την οποία ακτινοβολούν ενέργεια και μια αριθμητική τιμή που ονομάζεται σταθερά Planck.

Η Κβαντική χημεία περιλαμβάνει βαθιά χρήση πειραματικών και θεωρητικών μεθόδων:
\begin{itemize}
\item Οι πειραματικοί κβαντικοί χημικοί βασίζονται στην φασματοσκοπία, μέσω της οποίας η πληροφορία σχετικά με την κβάντιση της ενέργειας σε μοριακή κλίμακα μπορεί να αποκτηθεί. Τυπικές μέθοδοι είναι ή υπέρυθρη φασματοσκοπία και η πυρηνική μαγνητική αντήχηση (NMR) φασματοσκοπίας
\item Θεωρητική κβαντική χημεία, το έργο της οποίας βρίσκεται κάτω απο την κατηγορία της υπολογιστικής χημείας, στοχεύει να υπολογίσει τις προβλέψεις της κβαντικής θεωρίας καθώς τα άτομα και τα μόρια μπορούν να έχουν μόνο διακριτές ενέργειες. Λόγω του ότι αυτή η διαδικασία, όταν εφαρμοστεί σε πολύ-ατομικά είδη δημιουργεί υπολογιστικά προβλήματα, αυτοί οι υπολογισμοί γίνονται με την χρήση υπολογιστών αντί για την αναλυτική μέθοδο, και άλλες παραδοσιακές μεθόδους.
\end{itemize}
Κατά αυτούς τους τρόπους, η κβαντική χημεία μελετάει τα χημικά φαινόμενα. Στις αντιδράσεις, η κβαντική χημεία μελετάει την σταθερή κατάσταση των ατόμων και των μορίων, τις μη σταθερές καταστάσεις, και την μετάβαση των καταστάσεων ανάμεσα στις χημικές αντιδράσεις. Στους υπολογισμούς, η μελέτη της κβαντικής χημείας περιλαμβάνει χρήση ημι-εμπειρικών και άλλων μεθόδων που βασίζονται στις αρχές της κβαντικής μηχανικής, και αντιμετωπίζουν προβλήματα σχετικά με τον χρόνο. Πολλές μελέτες κβαντικής χημείας υποθέτουν ότι οι πυρήνες βρίσκονται σε κατάσταση ηρεμίας και πολλοί υπολογισμοί περιέχουν επαναληπτικές μεθόδους. Οι στόχοι της κβαντικής χημείας περιλαμβάνουν την αύξηση της ακρίβειας των αποτελεσμάτων για μικρά μοριακά συστήματα, και αύξηση του μεγέθους των μορίων που μπορούν να επεξεργαστούν. Ο χρόνος υπολογισμού αυξάνεται όσο αυξάνεται ο αριθμός των ατόμων.


\subsection{Προγράμματα}

\begin{apptable}{Κβαντική Χημεία}{chemistry}
Abinit & Επιτρέπει τον υπολογισμό της ολικής ενέργειας, πυκνότητα φόρτισης και ηλεκτρονική δομή του συστήματος αποτελούμενο από ηλεκτρόνια και πυρήνες μέσα στο DFT  & 1.3-2.7x & ΝΑΙ \\ \hline
ADF & Λογισμικό Θεωρίας Λειτουργικής Πυκνότητας (DFT) που επιτρέπει υπολογισμούς ηλεκτρονικών δομών πρώτης αρχής & 1.5-2x & ΝΑΙ \\ \hline
BigDFT & Εκτελεί θεωρία  λειτουργικής πυκνότητας λύνοντας τις εξισώσεις Kohn-Sham περιγράφοντας τα ηλεκτρόνια σε ένα υλικό & 2-25x & ΝΑΙ \\ \hline
CP2K & Πρόγραμμα που εκτελεί ατομιστικές και μοριακές εξομοιώσεις στερεών καταστάσεων, υγρών, μοριακών και βιολογικών συστημάτων & 2-7x & ΝΑΙ \\ \hline
GAMESS-UK & The general purpose ab initio molecular electronic structure program for performing SCF-, DFT- and MCSCF-gradient calculations & 8x & ΝΑΙ \\ \hline
GAMESS-US & Computational chemistry suite used to simulate atomic and molecular electronic structure & 1.3-2.9x & ΝΑΙ \\ \hline
Gaussian (In Development) & Predicts energies, molecular structures,and vibrational frequencies of molecular systems & - & ΝΑΙ \\ \hline 
GPAW & Real-space grid DFT code written in C and Python & 8x  & ΝΑΙ \\ \hline
LATTE & Density matrix computations & - & ΝΑΙ \\ \hline
MOLCAS & Methods for calculating general electronic structures in molecular systems in both ground and excited states & 1.1x & ΝΑΙ \\ \hline
MOPAC2013 & Semiempirical Quantum Chemistry & 2x & ΟΧΙ \\ \hline
NWChem & Calculations & 3-10x & ΝΑΙ \\ \hline
Octopus & Used for ab initio virtual experimentation and quantum chemistry calculations & 1.5-8x & - \\ \hline
Q-CHEM & Computational chemistry package designed for HPC clusters & 8-14x & - \\ \hline
QUICK & QUICK is a GPU-enabled ab intio quantum chemistry software package & 10-100x  & ΝΑΙ \\ \hline
TeraChem & Quantum chemistry software designed to run on NVIDIA GPU & 44-650x & ΝΑΙ \\ \hline
\end{apptable}

\section{Ψυχαγωγία}
\subsection{Εισαγωγή}
Μια μηχανή φυσικής είναι ένα πρόγραμμα υπολογιστή που παρέχει εξομοίωση συγκεκριμένων συστημάτων φυσικής, όπως δυναμική άκαμπτων σωμάτων, ανίχνευση σύγκρουσης, δυναμική υγρών, για χρήση σε πεδία όπως τα γραφικά υπολογιστών, παιχνίδια, κινούμενα σχέδια, ταινίες. Μια από τις κύριες χρήσεις τους είναι στα παιχνίδια υπολογιστών, στην οποία περίπτωση η εξομοίωση γίνεται σε πραγματικό χρόνο. O όρος χρησιμοποιείται γενικότερα για να περιγράψει οποιοδήποτε σύστημα λογισμικού που εξομοιώνει φυσικά φαινόμενα, όπως επιστημονικές εξομοιώσεις υψηλής απόδοσης.

Οι μηχανές φυσικής έχουν χρησιμοποιηθεί αρκετά στους υπερ-υπολογιστές από την δεκαετία του '80 για να εκτελέσουν μοντελοποίηση δυναμικών υγρών, όπου αναθέτουμε διανύσματα ισχύος σε σωματίδια, για να δείξουμε την κυκλοφορία. Λόγω των υψηλών απαιτήσεων σε ταχύτητα και ακρίβεια, ειδικοί επεξεργαστές δημιουργήθηκαν που είναι γνωστοί ως επεξεργαστές διανυσμάτων για να επιταχύνουν τους υπολογισμούς. Οι τεχνικές μπορούν να χρησιμοποιηθούν για να μοντελοποιήσουν πρότυπα καιρού για την πρόβλεψη καιρού, δεδομένα σήραγγας αέρα για σχεδιασμό αεροπλάνων και υποβρυχίων, και ανάλυση θερμικής απόδοσης για καλύτερο σχεδιασμό ψηκτρών για επεξεργαστές. Φυσικά μεγάλο ρόλο παίζει η ακρίβεια των υπολογισμών, αφού μικρές αποκλίσεις μπορούν να αλλάξουν δραστικά τα αποτελέσματα των υπολογισμών. Οι κατασκευαστές ελαστικών χρησιμοποιούν εξομοιώσεις φυσικής για να μελετήσουν πώς οι καινούριοι τύποι ελαστικών θα αποδίδουν σε συνθήκες βρεγμένου και στεγνού οδοστρώματος, χρησιμοποιώντας καινούρια υλικά και κάτω από διαφορετικές συνθήκες βάρους.

Υπάρχουν γενικά δύο τύποι μηχανών φυσικής. Οι πραγματικού χρόνου, και οι υψηλής ακρίβειας. Οι υψηλής ακρίβειας απαιτούν περισσότερη υπολογιστική δύναμη για να υπολογίσουν φυσικά φαινόμενα με ακρίβεια και χρησιμοποιούνται συνήθως από επιστήμονες αλλά και σε κινούμενα σχέδια. Οι πραγματικού χρόνου - χρησιμοποιούνται σε παιχνίδια υπολογιστών και σε άλλες μορφές διαδραστικού υπολογισμού - χρησιμοποιούν απλοποιημένους υπολογισμούς με μειωμένη ακρίβεια ώστε να επιτρέπουν στο παιχνίδι να αντιδράει σε αποδεκτό ρυθμό για την εμπειρία χρήσης.


\subsection{Παιχνίδια}
\subsubsection{Γενικά}
Στα περισσότερα παιχνίδια, η ταχύτητα των επεξεργαστών και η εμπειρία χρήσης είναι πιο σημαντικά από την ακρίβεια της εξομοίωσης. Αυτό μας οδηγεί σε σχεδιασμούς μηχανών φυσικής που παράγουν αποτελέσματα σε πραγματικό χρόνο αλλά αντιγράφουν φυσικά φαινόμενα μόνο για απλές περιπτώσεις. Τις περισσότερες φορές, η εξομοίωση είναι σχεδιασμένη να παρέχει μια φαινομενικά σωστή εκτίμηση, παρά απόλυτη ακρίβεια. Όμως μερικές μηχανές, απαιτούν μεγαλύτερη ακρίβεια σε σκηνές μάχης ή σε παιχνίδια τύπου παζλ. Οι κινήσεις χαρακτήρων στο παρελθόν χρησιμοποιούσαν φυσική άκαμπτων σωμάτων γιατί είναι γρηγορότερο και πιο εύκολο να υπολογιστεί, όμως τα τελευταία χρόνια τα παιχνίδια και οι ταινίες έχουν αρχίσει να χρησιμοποιούν φυσική μαλακών σωμάτων. Αυτού του τύπου η εξομοίωση χρησιμοποιείται επίσης για εφέ σωματιδίων, κίνηση υγρών και υφασμάτων. Μια μορφή εξομοίωσης δυναμικής υγρών χρησιμοποιείται για να εξομοιώσει νερό και άλλα υγρά αλλά και την ροή της φωτιάς και του καπνού στον αέρα.

\begin{figure}[h]
\centering
\includegraphics[width=0.5\linewidth]{havok-logo}
\caption{Λογότυπο της μηχανή φυσικής havok}
\end{figure}

\subsubsection{Ανίχνευση σύγκρουσης}
Η ανίχνευση σύγκρουσης συνήθως αναφέρεται στο υπολογιστικό πρόβλημα της ανίχνευσης της διασταύρωσης δύο η περισσότερων αντικειμένων. Αν και το θέμα έχει σχέση περισσότερο με την χρήση του στα παιχνίδια και σε άλλες εξομοιώσεις φυσικής, έχει και χρήσεις στην ρομποτική. Εκτός από την ανίχνευση του αν δύο αντικείμενα έχουν συγκρουστεί, τα συστήματα ανίχνευσης μπορούν να υπολογίσουν τον χρόνο της σύγκρουσης(Time Of Impact), και να αναφέρουν ένα σύνολο από σημεία διασταύρωσης. Η αντίδραση σύγκρουσης είναι η εξομοίωση του τι συμβαίνει όταν ανιχνευθεί μια σύγκρουση.
\subsection{Υλοποιήσεις}
\subsubsection{Μονάδα Επεξεργασίας Φυσικής}
Η μονάδα επεξεργασίας φυσικής (PPU) είναι ένας μικροεπεξεργαστής αποκλειστικά σχεδιασμένος για να χειρίζεται τους υπολογισμούς φυσικής, ειδικά σε μηχανές φυσικής των παιχνιδιών υπολογιστών. Η ιδέα είναι ότι αυτοί οι ειδικοί επεξεργαστές ελαφρύνουν το φόρτο εργασίας των CPUs, όπως μια κάρτα γραφικών εκτελεί υπολογισμούς γραφικών. Ο όρος αρχικά δημιουργήθηκε από την εταιρία Ageia για να περιγράψει τους επεξεργαστές PhysX στους καταναλωτές.
H NVIDIA απέκτησε την Ageia Technologies το 2008 και συνεχίζει να αναπτύσσει την πλατφόρμα PhysX και στο υλικό αλλά και στο λογισμικό. Απο την έκδοση 2.8.3, η υποστήριξη για κάρτες PPU σταμάτησε, και δεν κατασκευάζονται πλέον.  
\subsubsection{Υπολογισμοί γενικής χρήσης σε GPUs}
Η επιτάχυνση υλικού για υπολογισμούς φυσικής χρησιμοποιείται πλέον από τις GPUs που υποστηρίζουν υπολογισμό γενικής χρήσης. Η εκτέλεση φυσικών υπολογισμών σε GPUs είναι συνήθως αρκετά πιο γρήγορη απο ότι σε μια CPU, έτσι η απόδοση των παιχνιδιών βελτιώνεται και ή ροή εικόνας μπορεί να είναι πολύ πιο γρήγορη. Όμως η χρήση υπολογισμών φυσικής σε ένα παιχνίδι δημιουργεί επιπλέον φόρτο στην GPU. Έτσι, η χρήση ξεχωριστής μονάδας επεξεργασίας γραφικών για εκτελέσεις υπολογισμών φυσικής μπορεί να αποδώσει τα βέλτιστα αποτελέσματα. Το PhysX εκτελείται γρήγορα και αποδίδει μεγαλύτερο ρεαλισμό όταν εκτελείται στην GPU, αποφέροντας 10-20 φορές περισσότερα εφέ και οπτική πιστότητα απο ότι οι υπολογισμοί φυσικής που εκτελούνται σε μια κεντρική μονάδα επεξεργασίας τελευταίας τεχνολογίας. Το PhysX χρησιμοποιεί ετερογενή υπολογισμό για να αποδώσει την καλύτερη εμπειρία χρήσης. Καθώς το παιχνίδι εκτελείται, το σύστημα PhysX εκτελεί μέρη της τεχνολογίας στην CPU αλλά και άλλα μέρη στην GPU. Αυτό γίνεται ώστε να χρησιμοποιείται αποδοτικά το υλικό του υπολογιστή ώστε να παρέχουν την καλύτερη δυνατή εμπειρία στον χρήστη. Το πιο σημαντικό, είναι οτι η τεχνολογία PhysX μπορεί να κλιμακώνεται με την χρήση GPU, σε αντίθεση με άλλες ανταγωνιστικές υλοποιήσεις φυσικής.
\subsubsection{Σύγκριση}
Οι πιο σημαντικές μηχανές γραφικών που χρησιμοποιούνται σήμερα είναι οι παρακάτω:
\begin{itemize}
\item PhysX - Είναι μια μηχανή φυσικής πραγματικού χρόνου, από την NVIDIA. Είναι κλειστού κώδικα και χρησιμοποιείται σε πολλά παιχνίδια υπολογιστών και κονσολών. Υποστηρίζει μεγάλο αριθμό συσκευών.
\item Havok - Είναι μια μηχανή φυσικής που χρησιμοποιείται σε πολλά παιχνίδια υπολογιστών. 
\item ODE - Είναι μια μηχανή φυσικής που υποστηρίζει κυρίως ανίχνευση συγκρούσεων, δυναμικές άκαμπτων σωμάτων. Αποτελεί ανοιχτό και ελεύθερο λογισμικό. Έχει χρησιμοποιηθεί σε πολλά παιχνίδια και εφαρμογές. Αποτελεί δημοφιλής επιλογή για εφαρμογές εξομοίωσης ρομποτικής.
\item Newton Game Dynamics - Είναι μια μηχανή φυσικής ανοιχτού κώδικα που εξομοιώνει άκαμπτα σώματα σε παιχνίδια και άλλες εφαρμογές πραγματικού χρόνου. 
\item Bullet - Είναι μια μηχανή φυσικής που εξομοιώνει ανίχνευση συγκρούσεων, δυναμική άκαμπτων και μαλακών σωμάτων. Χρησιμοποιείται σε παιχνίδια υπολογιστών αλλά και για οπτικά εφέ σε ταινίες. Η βιβλιοθήκη της bullet physics είναι ελεύθερη και ανοιχτού κώδικα κάτω από την άδεια zlib.
\end{itemize}
\section{Μετατροπή}
\subsection{Εισαγωγή}
Η μετατροπή είναι η άμεση μετατροπή ψηφιακής κωδικοποίησης, όπως για παράδειγμα σε αρχεία μιας ταινίας(π.χ .mp4, avi), αρχεία ήχου(.wav .mp3), ή κωδικοποίηση χαρακτήρων. (π.χ UTF-8, ISO/IEC 8859). Αυτή γίνεται σε περιπτώσεις που η συσκευή στόχος δεν υποστηρίζει την μορφή του αρχείου ή χρειαζόμαστε μικρότερο μέγεθος αρχείου, ή για να μετατρέψουμε αρχεία παλαιότερης μορφής σε μια πιο σύγχρονη για καλύτερη υποστήριξη σε μελλοντικές εφαρμογές. 

Η μετατροπή χρησιμοποιείται συχνά στα λογισμικά προβολής βίντεο για να ελαττώσουμε το μέγεθος του αρχείου βίντεο. Μια διαδικασία που γίνεται συχνά είναι η μετατροπή από αρχεία MPEG-2(DVD) σε αρχεία μορφής MPEG-4, που ενσωματώνει σύγχρονους αλγόριθμους για καλύτερη ποιότητα εικόνας σε συνδυασμό με μικρότερο μέγεθος αρχείου.

\begin{figure}[h]
\centering
\includegraphics[width=0.5\linewidth]{h264-logo}
\caption{H264, το πιο δημοφιλές πρότυπο συμπίεσης βίντεο}
\end{figure}
\subsection{Μειονεκτήματα} 
Το μεγαλύτερο μειονέκτημα της μετατροπής σε απωλεστικές μορφές αρχείου είναι η μειωμένη ποιότητα. Τα τεχνουργήματα συμπίεσης συσσωρεύονται, οπότε κάθε διαδικασία μετατροπής δημιουργεί μια βαθμιαία απώλεια ποιότητας, που είναι γνωστή ως ψηφιακή απώλεια. Για αυτόν τον λόγο, η μετατροπή συνήθως δεν συνίσταται εκτός και αν δεν μπορούμε να την αποφύγουμε.

\subsection{Βίντεο}
H συμπίεση βίντεο χρησιμοποιεί σύγχρονες τεχνικές για να μειώσει πλεονασμούς στα δεδομένα βίντεο. Οι περισσότεροι αλγόριθμοι συμπίεσης βίντεο συνδυάζουν συμπίεση εικόνας και προσωρινή αποζημίωση κίνησης. Ο ήχος κωδικοποιείται παράλληλα με διαφορετικούς αλγόριθμους συμπίεσης άλλα συνήθως συνδυάζεται σε ένα πακέτο όταν ολοκληρωθεί η διαδικασία.

Οι περισσότεροι αλγόριθμοι συμπίεσης βίντεο χρησιμοποιούν απωλεστική συμπίεση, καθώς το ασυμπίεστο βίντεο απαιτεί πολύ μεγάλους ρυθμούς μετάδοσης. Αν και οι περισσότεροι αλγόριθμοι έχουν ένα παράγοντα συμπίεσης 3, μια τυπική συμπίεση βίντεο MPEG-4 μπορεί να έχει παράγοντα συμπίεσης από 20 έως και 200. Όπως σε όλες τις απωλεστικές διαδικασίες συμπίεσης, υπάρχει ένα δίλημμα μεταξύ της ποιότητας του βίντεο, το κόστος της επεξεργασίας της συμπίεσης και της αποσυμπίεσης, και των απαιτήσεων του συστήματος. Υπερβολικά συμπιεσμένο βίντεο μπορεί να δημιουργήσει οπτικά τεχνουργήματα.

Συνήθως ο τρόπος με τον οποίο λειτουργεί ένας αλγόριθμος συμπίεσης βίντεο είναι με ομάδες γειτονικών εικονοστοιχείων(pixel), που ονομάζονται blocks. Αυτές οι ομάδες από pixels, συγκρίνονται μεταξύ τους από μια εικόνα στην επόμενη, και ο αλγόριθμος αποστέλλει μόνο τις αλλαγές μεταξύ αυτών των block. Σε περιοχές του βίντεο που υπάρχει μεγαλύτερη κίνηση, ο αλγόριθμος πρέπει να συμπιέσει περισσότερα δεδομένα για να προλάβει τον μεγαλύτερο αριθμό εικονοστοιχείων που αλλάζουν. Συνήθως σε σκηνές με φωτιά, εκρήξεις, καπνούς, το αποτέλεσμα της συμπίεσης έχει μεγαλύτερη απώλεια ποιότητας, η αύξηση του ρυθμού μετάδοσης.

Μερικοί από τους σύγχρονους αλγόριθμους μετατροπής είναι οι παρακάτω
\begin{itemize}
\item Lagarith - Είναι ένας μη απωλεστικός αλγόριθμος ανοιχτού κώδικα που δίνει έμφαση στην ταχύτητα, στην υποστήριξη διαφόρων χώρου χρωμάτων(YV12,RGB,YUY2). Είναι ιδανικός για επεξεργασία και αποθήκευση αρχείων. Αν και υπάρχουν κάποιοι καλύτεροι μη απωλεστικοί αλγόριθμοι συμπίεσης, ο lagarith είναι ο πιο γρήγορος και έτσι έχει κερδίσει την υποστήριξη της κοινότητας.
\item VP9 - Είναι ένα ελεύθερο πρότυπο ανοιχτού κώδικα που αναπτύσσεται από την Google. Είναι ο διάδοχος του VP8. Σκοπός του είναι να μειώσει περισσότερο τον χώρο που απαιτείται από το βίντεο διατηρώντας την ίδια ποιότητα. 
\item H.264 - Είναι ένα πρότυπο συμπίεσης βίντεο, που είναι ίσως το πιο επιτυχημένο για την καταγραφή, συμπίεση, και αναμετάδοση περιεχομένου βίντεο. Χρησιμοποιείται για αναμετάδοση ψηφιακού σήματος τηλεόρασης, δορυφόρων, και από υπηρεσίες αναμετάδοσης στο διαδίκτυο.
\item H.265 - Είναι ένα πρότυπο συμπίεσης βίντεο, διάδοχος του επιτυχημένου H.264. Σκοπός του είναι ο διπλασιασμός της συμπίεσης διατηρώντας την ίδια ποιότητα. Υποστηρίζει αναλύσεις έως και 8192x4320.
\item Daala - Είναι μια τεχνολογία συμπίεσης από το ίδρυμα Xiph.Org. Χρησιμοποιεί περιτύλιξη μεταμορφώσεων για να μειώσει τα οπτικά τεχνουργήματα. Ο στόχος είναι η απόδοση του να ξεπεράσει τις δυνατότητες του VP9 και του H.265. 
\end{itemize}

\subsection{Συμπίεση δεδομένων γενετικής}
Οι γενετικοί αλγόριθμοι συμπίεσης είναι η τελευταία γενιά μη απωλεστικών αλγορίθμων για συμπίεση δεδομένων (συνήθως αλληλουχίες nucleotides) χρησιμοποιώντας συμβατικούς αλγορίθμους συμπίεσης αλλά και γενετικούς αλγόριθμους βελτιστοποιημένους στο συγκεκριμένο τύπο δεδομένων. Το 2012, μια ομάδα επιστημόνων απο το John Kopkins University ανακοίνωσε έναν αλγόριθμο συμπίεσης γενετικής(HapZipper), ο οποίος χρησιμοποιεί HapMap δεδομένα για να πετύχει συμπίεση 95\% μείωση στο μέγεθος αρχείου, πετυχαίνοντας καλύτερη συμπίεση σε πολύ καλύτερο χρόνο από ότι οι γνωστοί μη απωλεστικοί αλγόριθμοι. Άλλοι γενετικοί αλγόριθμοι συμπίεσης π.χ GenomeZip πετυχαίνουν μεγαλύτερη συμπίεση καταφέρνοντας αποθήκευση 6 δισεκατομμυρίων ανθρώπινων ζευγαριών γονιδιώματος σε 2.5 megabyte.
\section{Γεωπληροφορική}

\subsection{Σεισμική δραστηριότητα}
\subsubsection{Εισαγωγή}
Η ενσωματωμένη εξομοίωση σεισμικής δραστηριότητας είναι μια ανώτερη τεχνολογία για προσδιορισμό και απεικόνιση της δομικής βλάβης που δημιουργείται σε ένα σενάριο σεισμικής δραστηριότητας. Στο IES, όλα τα κτήρια που βρίσκονται στην περιοχή δοκιμής απεικονίζονται σαν δομικά μοντέλα, και η δομική ζημιά μπορεί να προσδιοριστεί με γραμμική και μη-γραμμική δυναμική δομική ανάλυση. Επειδή υπάρχουν πολλά κτήρια σε μια περιοχή δοκιμής, και κάθε κτήριο έχει διαφορετικά χαρακτηριστικά, ο υπολογισμός υψηλής ανάλυσης είναι απαραίτητος ώστε να εκτελεστεί η ανάλυση σε σχετικά μικρό χρόνο.\cite{geoinformatics-1}

\begin{figure}[h]
\centering
\includegraphics[scale=0.50]{ies}
\caption{Ενσωματωμένη εξομοίωση σεισμικής δραστηριότητας}
\end{figure}

Ένας καλός υποψήφιος για τον IES, είναι το OBASAN (Object Based Structural Analysis). Σε μια εξομοίωση σεισμικής δραστηριότητας, κάθε κτήριο σε μια περιοχή δοκιμών μπορεί να παρασταθεί ως δομικό μοντέλο για ανάλυση του κάθε αντικειμένου. Το γεωγραφικό πληροφοριακό σύστημα παρέχει δεδομένα δομικών σχημάτων στην μορφή των στοιχείων και πολυγώνων, και άλλα συστήματα βάσεων δεδομένων παρέχουν τύπους κτηρίων (π.χ δομή οπλισμένου σκυροδέματος, δομή χάλυβα, ξύλινη δομή, κ.α) Βασισμένο στον τρόπο σχεδιασμού της Ιαπωνίας, η ανάλυση μοντέλων των κτηρίων στον IES δημιουργούνται αυτόματα από τις διαστάσεις και τον τύπο του κάθε κτηρίου, και τότε το OBASAN μπορεί να χρησιμοποιηθεί για να αναλύσει αυτα τα μοντέλα δοκιμής ανάλυσης. Η δοκιμή ζημιά σε όλα τα κτήρια μιας περιοχής δοκιμής μπορεί να προβλεφθεί και η ολική ζημιά σε μια περιοχή πόλης μπορεί να απεικονιστεί από αυτήν την εξομοίωση σεισμού.\cite{geoinformatics-2}\cite{geoinformatics-3}

Τα πιο πολύπλοκα δομικά μοντέλα μπορούν να αυξήσουν την αξιοπιστία του προσδιορισμού δομικής βλάβης, π.χ τα μοντέλα μπορούν να χρησιμοποιηθούν για δυναμική ανάλυση δομικής απόκρισης. Μοντέλα δομικών αντικειμένων, μπορούν να χρησιμοποιηθούν για να δώσουν ένα πιο αξιόπιστο αποτέλεσμα. 

\subsubsection{OBASAN}
To OBASAN, είναι μια τεχνολογία που δημιουργήθηκε για την ανάλυση των δομικών μοντέλων ώστε να δέχεται τα αποτελέσματα των αντικειμένων (αξονική δύναμη, παραμόρφωση, καταπόνηση) και τα αποτελέσματα των κόμβων (π.χ ταχύτητα, επιτάχυνση, ισχύς) κάθε κτηρίου σε μια εξομοίωση σεισμικής δραστηριότητας, και έτσι το OBASAN μπορεί να παρέχει διάφορους τύπους αποτελεσμάτων δομικής ανάλυσης. Υπάρχουν δυνατότητες για χρήση υπολογισμού υψηλής απόδοσης ώστε να αυξηθεί η αποδοτικότητα και η δυνατότητα του συστήματος εξομοίωσης σεισμών (IES). 

\begin{figure}[h]
\centering
\hspace*{-1.0in}
\includegraphics[scale=0.50]{obasan}
\caption{OBASAN}
\end{figure}

Το OBASAN αναπτύχθηκε με αντικειμενοστραφή προγραμματισμό σε γλώσσα C++. Η γλώσσα C++ μπορεί να επεκταθεί ώστε να χρησιμοποιεί παράλληλο προγραμματισμό στην GPU. Όπως φαίνεται στο διάγραμμα, η αρχιτεκτονική του OBASAN έχει οργανωθεί συστηματικά και είναι εύκολο να κατανοηθεί. Λόγω του ότι υπάρχουν πολλά κτήρια σε μια εξομοίωση σεισμικής δραστηριότητας σε μια περιοχή δοκιμών, ο IES χρησιμοποιεί την αρχιτεκτονική OpenMPI για να ενεργοποιήσει παράλληλο υπολογισμό στην κεντρική μονάδα επεξεργασίας (CPU). Στην CPU, ένα υποσύνολο (π.χ 10) κτηρίων μπορεί να εκτελεστεί με τον μέγιστο αριθμό νημάτων σε έναν παράλληλο υπολογισμό. Όπως και στον IES όπου η αρχιτεκτονική OpenMPI χρησιμοποιείται για τον παράλληλο υπολογισμό των κτηρίων σε μια περιοχή δοκιμών, στο OBASAN χρησιμοποιείται η αρχιτεκτονική CUDA για την παράλληλη επεξεργασία της δυναμικής ανάλυσης δομικής απόκρισης(DSRA).

Λόγω του ότι η DSRA έχει να κάνει με πιο πολύπλοκα μοντέλα, η δομική ανάλυση των μοντέλων πρέπει να παρέχει μεγάλο αριθμό παραμέτρων, που εκφράζονται με μαθηματικές εξισώσεις που περιέχουν μεγάλους πίνακες. Για ένα τόσο προχωρημένο υπολογιστικό έργο, η τεχνική HPC είναι η συνηθέστερη προσέγγιση για την επίλυση των μαθηματικών εξισώσεων σε μικρό χρόνο. Στο HPC, η τεχνολογία GPGPU είναι σχεδιασμένη ώστε να αντιμετωπίζει μεγάλο αριθμό δεδομένων και μεγάλο αριθμό υπολογισμών σε μαθηματικές εξισώσεις. Στο GPGPU, η αρχιτεκτονική CUDA μπορεί να επιλεχτεί για να λύσει αυτές τις μαθηματικές εξισώσεις της DSRA. Το κάθε βήμα της ανάλυσης της DSRA, απαιτεί πολλαπλασιασμούς πινάκων για την δομική ανάλυση.

Στο παρακάτω σχήμα μπορούμε να δούμε ένα παράδειγμα χρόνου εκτέλεσης πολλαπλασιασμού πινάκων. Οι πίνακες Α,Β και C υποθέτουμε ότι είναι τετράγωνοι πίνακες και το μέγεθος τους είναι Ν * Ν. Ο κώδικας C++ του πολλαπλασιασμού πινάκων εκτελείται σε μια CPU (Intel Xeon), ενώ οι εντολές CUDA σε μια GPU(Tesla M2050). 

\begin{figure}[h]
\centering
\includegraphics[scale=0.50]{matrix-multi}
\caption{GPU and CPU matrix multiplication times}
\end{figure}

Στην DRSA, το μέγεθος πίνακα ορίζεται απο τον συνολικό αριθμό των μοιρών της ελευθερίας (DOF) σε ένα κτήριο, οπότε το μέγεθος του πίνακα εξαρτάται απο τις διαστάσεις του κάθε κτηρίου. Τα μικρότερα κτήρια όπως τα σπίτια αναπαρίστανται απο μικρότερους πίνακες και τα μεγαλύτερα κτήρια αναπαρίστανται απο μεγαλύτερους πίνακες. Όπως μπορούμε να δούμε στον πίνακα, οι GPUs είναι πολύ πιο γρήγορες απο τις CPUs σε κάθε περίπτωση πολλαπλασιασμού πινάκων. Από τα δεδομένα του πίνακα, μπορούμε να απεικονίσουμε τα αποτελέσματα σε γράφημα. Στο παρακάτω γράφημα βλέπουμε την σχέση που έχει ο χρόνος εκτέλεσης με το μέγεθος του πίνακα, τόσο στην CPU όσο και στην GPU. Οι μεγαλύτεροι πίνακες χρειάζονται περισσότερο χρόνο για να εκτελέσουν τον υπολογισμό πινάκων. Όσο μεγαλύτερο το μέγεθος του πίνακα, τόσο μεγαλύτερο το πλεονέκτημα της GPU έναντι της CPU. Για ένα πίνακα 2000 * 2000 στοιχείων, οι GPUs είναι περίπου 1000 φορές γρηγορότερες απο τις CPUs. Για αυτό ο προγραμματισμός γενικού σκοπού είναι καλή επιλογή για την δυναμική ανάλυση δομικής απόκρισης.

\begin{figure}[h]
\centering
\includegraphics[scale=0.50]{graph-multi}
\caption{GPU and CPU matrix multiplication times}
\end{figure}

Τα παραπάνω αποτελέσματα, είναι μια απλή τεχνική CUDA. Ο προγραμματισμός υψηλής απόδοσης της DSRA μπορεί να αυξηθεί με την χρήση βιβλιοθηκών CUDA. Οι CUDA βιβλιοθήκες, π.χ CUBLAS(Γραμμική Άλγεβρα),CUFFT(Μετασχηματισμός Fast Fourier), μπορούν να λύσουν τις μαθηματικές εξισώσεις της DSRA. Γενικότερα, όλες οι βιβλιοθήκες CUDA είναι βελτιστοποιημένες για εκτέλεση υψηλής απόδοσης. Επιπλέον, το OBASAN με την χρήση CUDA μπορεί να επιτύχει καλή κλιμάκωση στον IES. 
\subsubsection{Συμπεράσματα}
Για να μειώσουν την κοινωνική και οικονομική επίπτωση των σεισμών, οι ερευνητές προσπαθούν να εφαρμόσουν καινούριες τεχνολογίες και τα τελευταία δεδομένα στον IES. Η αριθμητική εξομοίωση χρησιμοποιείται για να προσδιορίσει πιθανές ζημιές σε σενάρια σεισμών. Σε μια εξομοίωση σεισμών, η σωστή πρόβλεψη της ζημιάς στην κοινωνία είναι πολύ σημαντική στην διαχείριση κρίσεων. Αυτή η προσέγγιση στο μαθηματικό πρόβλημα με την χρήση τεχνολογιών προγραμματισμού υψηλής απόδοσης και CUDA, μπορεί να εφαρμοστεί σε διάφορα πεδία της μηχανικής. \cite{geoinformatics-4}\cite{geoinformatics-5}

Με την χρήση των προσδιορισμών, ένα πιο σταθερό πλάνο μπορεί να χρησιμοποιηθεί για κατασκευή καινούριων κτηρίων για να εμποδίσει δομικές ζημιές σε σενάρια σεισμών. Επιπλέον, σχέδια για εκκένωση κτηρίων και μεταφορά ατόμων μπορούν να δημιουργηθούν για να αναλύσουν τις εξόδους κινδύνου, κάτι που μπορεί να προστατέψει και να σώσει ανθρώπινες ζωές.

\subsection{Προγράμματα}

\begin{apptable}{Γεωπληροφορική}{geoinformatics}
DigitalGlobe Advanced Ortho Series & Geospatial Visualization & 50x & ΝΑΙ \\ \hline
Eternix Blaze Terra & Geospatial Visualization  & 50x & ΝΑΙ \\ \hline
Exelis (ITT) ENVI & Geospatial Visualization & 70x & ΝΑΙ \\ \hline
GAIA & advanced geospatial query capability,heatmap generation, and distributed rasterization services & - & ΝΑΙ \\ \hline
GeoWeb3d Desktop & Geospatial Visualization & - & ΝΑΙ \\ \hline
Incogna GIS  & Geospatial Visualization & 50x & ΝΑΙ \\ \hline
LuciadLightspeed & Geospatial Visualization and analysis & & ΝΑΙ \\ \hline
Manifold Systems & Full-featured GIS, vector/raster processing \& analysis & - & ΝΑΙ \\ \hline
MrGeo & Geospatial Visualization  & - & ΝΑΙ \\ \hline
OpCoast SNEAK  & Electromagnetic signals propagation modeling for complex urban and terrain environments & 100x & ΝΑΙ \\ \hline
PCI Geomatics GXL  & Geospatial Visualization & 20-60x & ΝΑΙ \\ \hline
\end{apptable}

\begin{apptable}{Κατασκευαστική Μηχανική}{structural}
Abaqus/Standard & Simulation and analysis tool for structural mechanics & & ΝΑΙ \\ \hline
ANSYS Mechanical & Simulation and analysis tool for structural mechanics & & ΝΑΙ \\ \hline
DS Exsight & Uses Abaqus Standard for GPU computing & & ΟΧΙ \\ \hline
DS DesignSight & Uses Abaqus Standard for GPU computing & & ΟΧΙ \\ \hline
Impetus Afea & Predicts large deformations of structures and components exposed to extreme loading conditions & & ΝΑΙ \\ \hline
LS-DYNA Implicit  & Multiphysics simulation package used & & ΝΑΙ \\ \hline
MSC Nastran & Simulation and analysis tool for structural mechanics & & ΝΑΙ \\ \hline
MSC Marc & Simulation and analysis tool for structural mechanics & & ΝΑΙ \\ \hline
NX Nastran & Simulation and analysis tool for structural mechanics & & ΝΑΙ \\ \hline
OptiStruct & Simulation and analysis tool for structural mechanics & & ΝΑΙ \\ \hline
PAM-CRASH Implicit & Multiphysics simulation package used & & ΟΧΙ \\ \hline
\end{apptable}

\begin{apptable}{Σεισμική δραστηριότητα}
Acceleware AxRTM AxKTM & Seismic Processing & & ΝΑΙ \\ \hline
CGGV GeoVation & Seismic Processing & & ΝΑΙ \\ \hline
CGG-Veritas Inside Earth & Seismic Interpretation & & ΝΑΙ \\ \hline
ffA Geoteric & Seismic Interpretation  & & ΝΑΙ \\ \hline
ffA SEA3D Pro  & Seismic Interpretation  & & ΝΑΙ \\ \hline
ffA SVI Pro & Seismic Interpretation  & & ΝΑΙ \\ \hline
GeoMage Multifocusing & Seismic imaging  & & ΝΑΙ \\ \hline
GeoStar Seismic Suite & Seismic Processing  & & ΝΑΙ \\ \hline
HUE Headwave Suite & Seismic Interpretation  & & ΝΑΙ \\ \hline
HUE HUEspace & Seismic Interpretation  & & ΝΑΙ \\ \hline
OpenGeo Solutions OpenSeis & Seismic Processing  & & ΝΑΙ \\ \hline
Paradigm Echos RTM & Seismic Processing  & & ΝΑΙ \\ \hline
Paradigm Geophysical VoxelGeo & & & ΝΑΙ \\ \hline
Paradigm SKUA & Reservoir Modeling  & & ΝΑΙ \\ \hline
Panorama Tech & Seismic Processing, Modeling & & ΝΑΙ \\ \hline
Roxar RMS & Reservoir Modeling  & & ΝΑΙ \\ \hline
Schlumberger Omega2 RTM & Seismic Processing  & & ΝΑΙ \\ \hline
Seismic City Prestack Interpretation & Seismic Processing  & & ΝΑΙ \\ \hline
SpectraSeis & Seismic Processing  & & ΝΑΙ \\ \hline
Stoneridge Technologies GAMPACK & Reservoir Simulation & & ΝΑΙ \\ \hline
Tsunami A2011 & Seismic Processing/Imaging package  & & ΝΑΙ \\ \hline
Tsunami RTM  & Seismic Processing  & & ΝΑΙ \\ \hline
\end{apptable}


\section{Αστροφυσική}
Οι πρώτες καταγραφές της θεωρητικής αστρονομίας χρονολογούνται το 1550-1292 π.Χ Οι υπολογισμοί που βρέθηκαν σχηματισμένοι σε αιγυπτιακούς τάφους δείχνουν ότι υπήρχαν τεχνικές για την αναγνώριση και την καταγραφή προτύπων στον ουρανό, ένα πεδίο της αστρονομίας που είναι χρήσιμο ακόμα και σήμερα για την καταγραφή και χαρτογράφηση.

Μπορεί οι σημερινοί ερευνητές να μην χρησιμοποιούν τα αρχαία εργαλεία, αλλά η ανάπτυξη τους έγινε σταδιακά. Το μέγεθος των τηλεσκοπίων χρειάστηκε 400 χρόνια για να μεγαλώσει από 1 τετραγωνικό μέτρο σε 110 τετραγωνικά μέτρα. Οι ψηφιακοί υπολογιστές εμφανίστηκαν στο τέλος του 1940 με υπολογιστική ταχύτητα περίπου 100 floating point λειτουργιών το δευτερόλεπτο (FLOPS), και εξελίχθηκαν σε περίπου $3*10^{16}$ FLOPS σε λιγότερο από 65 χρόνια. Αυτή η επανάσταση των υπολογιστών συνεχίζεται ακόμα και σήμερα, και έχει οδηγήσει σε ένα καινούριο κομμάτι έρευνας στο οποίο οι εγκαταστάσεις δεν βρίσκονται στο ψηλότερο βουνό του κόσμου, αλλά στο διπλανό δωμάτιο. Οι Αστρονόμοι κατανόησαν γρήγορα ότι μπορούν να χρησιμοποιήσουν τους υπολογιστές για να καταγράψουν, αναλύσουν, αρχειοθετήσουν, τις τεράστιες ποσότητες πληροφοριών που καταγράφονται από τις εκστρατείες παρατηρητών.\cite{astrophysics} \\

Την μεγαλύτερη όμως επίδραση στον τρόπο με τον οποίο αντιμετωπίζουν οι αστρονόμοι τα αναπάντητα ερωτήματα τους, έχει το πεδίο της εξομοίωσης. Με την χρήση των υπολογιστών, είναι δυνατόν να μελετήσουμε την λειτουργία του Διαγαλαξιακούς κενού, την φυσική των μαύρων τρυπών, κ.α


\begin{apptable}{Φυσική}{physics}
Chemora & Το Chemora είναι ένα σύστημα που εκτελεί εξομοιώσεις συστημάτων που περιγράφονται από διαφορικές εξισώσεις που εκτελούνται σε επιταχυνόμενα υπολογιστικά συμπλέγματα & & ΝΑΙ \\ \hline
Chroma & Δικτυωτή Κβαντική Χρωμοδυναμική (LQCD) & & ΝΑΙ \\ \hline
ENZO & Τρισδιάστατος AMR κώδικας μορφής block για σχηματισμό κοσμολογικής μορφής & & ΝΑΙ \\ \hline
GTC & Εξομοιώνει μικρο-διαταραχές και μεταφορές σε μαγνητικά περιορισμένη σύντηξη πλάσματος & & ΝΑΙ \\ \hline
GTS & Εξομοιώνει μικρο-διαταραχές και την κίνηση των φορτισμένων σωματιδίων και αλληλεπιδράσεις σε σύντηξη πλάσματος & & ΝΑΙ \\ \hline
MILC & Δικτυωτή Κβαντική Χρωμοδυναμική (LQCD) εξομοιώνει πως δομούνται τα σωματίδια στοιχείων και δεσμεύονται από την δυνατή ισχύ ώστε να δημιουργήσουν μεγαλύτερα σωματίδια όπως πρωτόνια και νετρόνια & & ΝΑΙ \\ \hline
PIConGPU & Ένα σχεσιακό σύστημα σωματιδίων σε κελί που περιγράφει την δυναμική ενός πλάσματος υπολογίζοντας την κίνηση των ηλεκτρονίων και ιόντων που περιγράφεται από την εξίσωση Maxwell-Vlasov. & & ΝΑΙ \\ \hline
QUDA & Βιβλιοθήκη για δικτυωτούς QCD υπολογισμούς με χρήση GPUs & & ΝΑΙ \\ \hline
RAMSES  & Εξομοιώνει προβλήματα αστροφυσικής σε διαφορετικές κλίμακες(π.χ σχηματισμούς αστερισμών, δυναμικές γαλαξίων, σχηματισμό κοσμολογικής μορφής) & & ΝΑΙ \\ \hline
\end{apptable}
\chapter{Προσδιορισμός του προβλήματος και διεθνής εμπειρία}
\chapter{Μεθοδολογία και Εφαρμογή}
\section{mpla}
Το πεδίο του προγραμματισμού γενικής χρήσης έχει αναπτυχθεί με ταχύτατους ρυθμούς τα τελευταία χρόνια, έτσι ώστε τώρα υπάρχουν αρκετές υλοποιήσεις για τον προγραμματισμό των μονάδων επεξεργασίας γραφικών. Πρόσφατα, έχουν γίνει προσπάθειες δημιουργίας προτύπων.

Ο προγραμματισμός των GPUs αναπτύχθηκε όταν το CUDA και το Stream κατέφθασαν στο τέλος του 2006. Αυτές οι διεπαφές και οι γλώσσες, σχεδιάστηκαν απο τις εταιρίες κατασκευής των GPUs σε πολύ κοντινή σχέση με το υλικό, το οποίο αποτέλεσε μεγάλο βήμα προς ένα πιο εύχρηστο, ταιριαστό και μελλοντικά-ασφαλές προγραμματιστικό μοντέλο. Η ανοιχτή γλώσσα προγραμματισμού (OpenCL) δημιουργήθηκε για να παρέχει ένα γενικό API ετερογενή υπολογισμού σε διάφορες μορφές παράλληλων συσκευών, συμπεριλαμβανομένου μονάδων επεξεργασίας γραφικών, πολυπύρηνων κεντρικών μονάδων επεξεργαστών, κ.α
%
%	Copyright 2014 George Sofianos
%
%	Licensed under the EUPL, Version 1.1 or – as soon they
%	will be approved by the European Commission - subsequent
%	versions of the EUPL (the "Licence");
%	You may not use this work except in compliance with the
%	Licence.
% 	You may obtain a copy of the Licence at:
%
%	http://ec.europa.eu/idabc/eupl
%
%	Unless required by applicable law or agreed to in
%	writing, software distributed under the Licence is
%	distributed on an "AS IS" basis,
%	WITHOUT WARRANTIES OR CONDITIONS OF ANY KIND, either
%	express or implied.
%	See the Licence for the specific language governing
%	permissions and limitations under the Licence.
%
\section{Άδεια χρήσης}
\subsection{Εισαγωγή}
Η άδεια χρήσης λογισμικού, είναι ένα νομικό μέσο (συνήθως της μορφής του δικαίου των συμβάσεων) το οποίο ορίζει την χρήση και την διαμοίραση του λογισμικού. Σύμφωνα με τον Ευρωπαϊκό νόμο, όλα τα προγράμματα υπολογιστών προστατεύονται από πνευματικά δικαιώματα. Μια τυπική άδεια χρήσης, παρέχει διάφορα δικαιώματα στον τελικό χρήστη για την χρήση του λογισμικού, που σε άλλη περίπτωση θα ήταν καταπάτηση πνευματικών δικαιωμάτων.
Συνήθως οι άδειες χρήσεις λογισμικού εντάσσονται σε μια από τις παρακάτω κατηγορίες, ενώ η διαφορά τους βρίσκεται στους όρους με τους οποίους ο τελικός χρήστης μπορεί στην συνέχεια να διαμοιράσει το λογισμικό.
\subsubsection{Άδειες αποκλειστικής ιδιοκτησίας}
Το σήμα κατατεθέν του λογισμικού αποκλειστικής ιδιοκτησίας είναι ότι ο εκδότης επιτρέπει την χρήση μιας η περισσότερων αντίτυπων του λογισμικού, κάτω από μια συμφωνητικού τύπου άδεια χρήσης (EULA), αλλά η ιδιοκτησία των αντίτυπων παραμένει στον εκδότη. Αυτό το χαρακτηριστικό αυτής της άδειας, σημαίνει ότι συγκεκριμένα δικαιώματα σχετικά με το λογισμικό έχουν παρακρατηθεί από τον εκδότη του λογισμικού. Γιαυτό είναι συνηθισμένο, στις EULA να υπάρχουν όροι που ορίζουν της χρήσεις του λογισμικού, όπως ο αριθμός των εγκαταστάσεων που επιτρέπονται από τους όρους της διανομής.\\
Η μεγαλύτερη επίδραση αυτού του τύπου άδειας, είναι ότι, αν η ιδιοκτησία του λογισμικού παραμένει στον εκδότη, τότε ο τελικός χρήστης πρέπει να αποδεχτεί την άδεια του λογισμικού, ενώ σε περίπτωση που δεν το κάνει, δεν μπορεί να χρησιμοποιήσει το λογισμικό. Σε αυτού του τύπου άδειες περιέχεται συνήθως και μια εκτενής λίστα με λειτουργίες που απαγορεύονται, όπως για παράδειγμα reverse engineering, ταυτόχρονη χρήση λογισμικού από πολλούς χρήστες, κ.α
\subsubsection{Άδειες ελεύθερου και ανοιχτού κώδικα}
Οι άδειες ελεύθερου και ανοιχτού κώδικα, κατατάσσονται σε μια από τις δύο κατηγορίες: 
\begin{itemize}
\item Permissive άδειες - Αυτές οι άδειες έχουν λιγότερες απαιτήσεις για το πώς θα χρησιμοποιηθεί και διανεμηθεί το λογισμικό. Παράδειγμα permissive άδειας είναι η άδεια BSD και η άδεια MIT, που δίνουν απόλυτη ελευθερία για χρήση,μελέτη, και ιδιωτική τροποποίηση του λογισμικού, και περιέχει μόνο ελάχιστες απαιτήσεις για την διανομή. Αυτό δίνει σε κάποιον την δυνατότητα να χρησιμοποιήσει τον κώδικα σαν μέρος ενός λογισμικού κλειστού κώδικα, ή λογισμικού αποκλειστικής ιδιοκτησίας.
\item Copyleft άδειες - Αυτές οι άδειες έχουν ως στόχο να διατηρήσουν τις ελευθερίες που δίνονται στον χρήστη. Ένα παράδειγμα copyleft άδειας χρήσης λογισμικού, είναι το GPL. Αυτή η άδεια, στοχεύει στο να δώσει στους χρήστες την ελευθερία για χρήση,μελέτη, και ιδιωτική τροποποίηση του λογισμικού, και αν ο χρήστης εμμείνει στους όρους και προϋποθέσεις του GPL, ελευθερία για διανομή του λογισμικού και κάθε τροποποίησης που έχει γίνει σε αυτό. Για παράδειγμα, κάθε τροποποίηση που γίνεται από τον χρήστη πρέπει να συνοδεύεται από τον πηγαίο κώδικα των αλλαγών όταν το λογισμικό διανεμηθεί, και η άδεια για οποιαδήποτε παράγωγο του έργου να μην έχει τους οποιεσδήποτε περιορισμούς εκτός από αυτούς που το GPL ορίζει.
\end{itemize}
\subsection{Επιλογή}
Μετά από εκτενή αναζήτηση και έρευνα πάνω στις άδειες χρήσεις λογισμικου, αποφασίστηκε να χρησιμοποιηθούν δύο άδειες για την εκπόνηση της πτυχιακής εργασίας. Μία για τον πηγαίο κώδικα και μία για το τελικό αποτέλεσμα. 
\subsubsection{Πηγαίος κώδικας}
Η άδεια χρήσης EUPL 1.1 κρίθηκε η καλύτερη επιλογή για την πτυχιακή εργασία. Ο λόγος είναι ότι

\subsubsection{Βιβλίο}
Για την επιλογή της άδειας χρησιμοποιήθηκε ο οδηγός στο site της creativecommons \url{http://creativecommons.org/choose/}
Επέλεξα την άδεια Attribution-ShareAlike 4.0 International \ccbysa γιατί θεωρώ ότι η πτυχιακή εργασία μου μπορεί να βοηθήσει στο μέλλον και άλλους φοιτητές να γνωρίσουν την τεχνολογία του \LaTeX ,αλλά ταυτόχρονα είναι και συμβατή με την άδεια χρήσης του πηγαίου κώδικα.
\vfill
Το περιεχόμενο του παρόντος κειμένου υπάγεται σε Άδεια Χρήσης Creative Commons Attribution 4.0 \url{http://creativecommons.org/licenses/by-sa/4.0/}
\subsection{Παρουσίαση}
\subsubsection{Blender και Bullet Physics}
Το Blender είναι ένα επαγγελματικό πρόγραμμα τρισδιάστατων γραφικών το οποίο είναι ελεύθερο και ανοιχτού κώδικα. Χρησιμοποιείται για την δημιουργία ταινιών, οπτικών εφέ, τέχνης, τρισδιάστατων μοντέλων εκτύπωσης, διαδραστικές τρισδιάστατες εφαρμογές, και παιχνίδια υπολογιστών. Τα χαρακτηριστικά του Blender περιέχουν τρισδιάστατη μοντελοποίηση, δημιουργία υφών, εξομοίωση ρευστών και καπνού, εξομοίωση μαλακών σωμάτων, ανταπόδοση σκηνών. Περιέχει επίσης και μηχανή κατασκευής παιχνιδιών.
\begin{figure}[h]
\centering
\includegraphics[scale=0.7]{blender}
\caption{Λογότυπο Blender 3D\cite{figure-24}}
\end{figure}

To Blender επιλέχτηκε για την παρουσίαση της πτυχιακής εργασίας, με εκπαιδευτικό παράδειγμα για το πως μπορεί να επιταχυνθεί η ανταπόδοση γραφικών και εξομοίωση φυσικής σε ένα παράδειγμα γραφικών.

H μηχανή εξομοίωσης φυσικής επιλέχτηκε για την παρουσίαση της πτυχιακής εργασίας. H Bullet είναι μια μηχανή φυσικής που εξομοιώνει ανίχνευση συγκρούσεων, δυναμική άκαμπτων και μαλακών σωμάτων. Χρησιμοποιείται σε παιχνίδια υπολογιστών αλλά και για οπτικά εφέ σε ταινίες. Η βιβλιοθήκη της bullet physics είναι ελεύθερη και ανοιχτού κώδικα κάτω από την άδεια zlib. Η επιλογή της μηχανής φυσικής Bullet έγινε λόγω της μεγάλης δημοτικότητας που έχει και της καταξιωμένης εφαρμογής της σε γνωστές ταινίες και παιχνίδια, συμπεριλαμβανομένου των: Toy Story 3, Grand Theft Auto IV,V, 2012(ταινία), Hancock, Megamind, και λόγω της ενσωμάτωση της στο πρόγραμμα γραφικών Blender.

\begin{figure}[h]
\centering
\includegraphics[width=0.7\linewidth]{bullet}
\caption{Λογότυπο μηχανής φυσικής Bullet Physics\cite{figure-25}}
\end{figure}

\subsubsection{Netbeans και JOCL}
Το Netbeans είναι ένα περιβάλλον ανάπτυξης λογισμικού που χρησιμοποιείται κυρίως σε συνδυασμό με την γλώσσα προγραμματισμού Java, αλλά και με άλλες γλώσσες, όπως PHP,C/C++, και HTML5. Αποτελεί επίσης μια πλατφόρμα ανάπτυξης για επιτραπέζιες εφαρμογές Java, κ.α. Το Netbeans είναι ανοιχτού κώδικα και υποστηρίζει έλεγχο εκδόσεων.

Ο λόγος της επιλογής του Netbeans είναι ότι υπάρχει υποστήριξη για Java Bindings για προγραμματισμό OpenCL μέσω της βιβλιοθήκης JOCL. To JOCL υποστηρίζει όλες τις υλοποιήσεις OpenCL που υπάρχουν αυτήν την στιγμή.
\begin{itemize}
\item AMD APP SDK με υποστήριξη για OpenCL 1.2
\item NVIDIA οδηγοί με υποστήριξη για OpenCL 1.1
\item OpenCL σε OSX
\item Intel OpenCL SDK
\end{itemize}

Το JOCL χρησιμοποιήθηκε για προγραμματισμό εκπαιδευτικού παραδείγματος με την χρήση OpenCL και μελετήθηκε η απόδοση διαφόρων παραδειγμάτων σε CPU και GPU αντίστοιχα. Καθώς η Java είναι μια γλώσσα που έχουμε διδαχτεί ενδελεχώς κατά την διάρκεια των σπουδών μας, θεωρώ ότι είναι χρήσιμη και ίσως η πιο εύκολη μετάβαση από τον παραδοσιακό προγραμματισμό σε προγραμματισμό GPU.
Αντίστοιχα bindings παρέχονται και για την τεχνολογία CUDA, μέσω της βιβλιοθήκης Jcuda.

\subsubsection{Sfera}
Το Sfera είναι ένα εκπαιδευτικό παιχνίδι με μεγάλη επιστημονική σημασία, ανεπτυγμένο από τον προγραμματιστή David Dade Bucciarelli. Το Sfera χρησιμοποιεί κυρίως OpenCL και OpenGL, και εκτελεί ανταπόδοση γραφικών με ιχνογράφηση ακτίνας πραγματικού χρόνου. Χρησιμοποιεί την μηχανή φυσικής Bullet. Το sfera βασίζεται στο πρόγραμμα SmallLiuxGPU v2.0 Μερικά από τα χαρακτηριστικά του είναι τα εξής:
\begin{itemize}
\item Πολλαπλές μηχανές ανταπόδοσης: μονο-νηματική, πολυ-νηματική, μία ή περισσότερες GPUs
\item Πολλαπλά υλικά που εκπέμπουν φως: καθρεύτες, γυαλί, μέταλλο, κράμα,ματ
\item Σχεδιασμός Υφών
\item Βάθος πεδίου
\item Φορητό: υλοποιήσεις σε Windows, Linux, MacOS, κ.α
\end{itemize}

\subsubsection{Compubench}
Το Compubench είναι ένα πρόγραμμα μετρήσεων απόδοσης OpenCL για τον έλεγχο και την σύγκριση της παράλληλης υπολογιστικής απόδοσης των CPUs,GPUs και επιταχυντές επιτραπέζιων και φορητών συσκευών. Περιλαμβάνει πολλές δοκιμές απόδοσης όπως αναγνώριση προσώπων, εξομοίωση επιφάνειας ωκεανών, εξομοίωση σωματιδίων, σύνθεση και μετατροπή βίντεο, εξόρυξη κρυπτο-νομισμάτων, κ.α.
\subsubsection{GPU Caps Viewer}
Το GPU Caps Viewer είναι ένα ελεύθερο και πλούσιο σε χαρακτηριστικά πρόγραμμα το οποίο περιγράφει τις δυνατότητες της GPU, συμπεριλαμβανομένου του τύπου της GPU, του μεγέθους της μνήμης, την έκδοση της υποστηριζόμενης έκδοσης OpenGL, ενώ περιέχει και διαδικασίες ελέγχου σταθερότητας και έντασης. Το GPU Caps Viewer περιέχει αρκετά παραδείγματα της χρήσης της τεχνολογίας OpenCL με χρήση γραφικών, ενώ βρίσκεται ανάμεσα στα πρώτα προγράμματα που υποστήριξαν την τεχνολογία OpenCL.
\subsubsection{Artifacts και demoscene}
Το παράδειγμα παρουσίασης Artifacts είναι μια παρουσίαση οπτικών εφέ και μουσικής, που κέρδισε την πρώτη θέση στον αντίστοιχο διαγωνισμό της demoscene στο TokyoDemoFest 2013. Είναι ένα καλό παράδειγμα των δυνατοτήτων της GPU τόσο από την οπτική γωνία των γραφικών όσο και από την οπτική γωνία της παρουσίασης των δυνατοτήτων της υπολογιστικής δύναμης των GPU.

dsadsa
\printbibliography[title={Βιβλιογραφία}]
\end{document}